\documentclass[10pt,a4paper]{article}

% ----------------- Paquetes -------------------%
\usepackage[T1]{fontenc}
\usepackage[utf8]{inputenc}
\usepackage[a4paper,margin=2.5cm]{geometry}
\usepackage{mathptmx}             
\usepackage{changepage} 
\usepackage{setspace}
\usepackage[spanish]{babel}
\usepackage[labelfont=bf,labelsep=space]{caption}
\usepackage{amssymb}
\usepackage{amsmath}
\usepackage{ebgaramond}
\usepackage{placeins}
\usepackage{etoolbox}
\usepackage{setspace}
\usepackage{parskip} 
\usepackage{tcolorbox}
\usepackage{needspace}
\usepackage{xparse}
\usepackage{ocgx2}
\usepackage{hyperref}
\usepackage{hypcap}



%%%%%%%%%%%%%%%%%%%%%%%%%%%%%%%%%%%%%%%%%%%%%%%%%%%%%%%%%%%%%%%%
% --- Fija primero tus valores globales "buenos" ---
\setstretch{1.2}                 % Interlineado global
\setlength{\parskip}{0.7\baselineskip}
\setlength{\parindent}{0pt}
\newlength{\GlobalParskip}
\setlength{\GlobalParskip}{\parskip}

\newcounter{eqnum}[subsection]
\renewcommand{\theeqnum}{\arabic{eqnum}}
\newcounter{alphaitem}
\newcounter{numitem}

%% ------------------- Recuadros----------------------%%
\newcommand{\puntosclave}[1]{%
  \begin{tcolorbox}[
    colframe=black,
    title={Puntos clave},
    before upper={\setlength{\parindent}{0.5cm}\setlength{\parskip}{0.5\baselineskip}}
  ]
  #1
  \end{tcolorbox}
}

\newcommand{\modificaciones}[1]{
  \begin{tcolorbox}[
    colback=white,
    colframe=red,
    title={Modificaciones a realizar},
    before upper={\setlength{\parindent}{0.5cm}\setlength{\parskip}{0.5\baselineskip}}
  ]
  #1
  \end{tcolorbox}
}

%% ------------------- Preguntas TEST----------------------%%
% Opciones: radio (single-choice)
\NewDocumentCommand{\OptRadio}{m m m}{%
  \noindent\CheckBox[radio,name=#1,value=#2,width=1.2ex,height=1.2ex]{}%
  \;\textbf{#2.}~#3\par
}

% Opciones: checkbox (multi-choice)
\NewDocumentCommand{\OptCheck}{m m}{%
  \noindent\CheckBox[name=#1,width=1.2ex,height=1.2ex,bordercolor={0 0 0}]{}%
  \;#2\par
}

% Pregunta single-choice (radio)
% Args: 1=id, 2=enunciado, 3=opciones, 4=solución+justificación, 5=imagen (o {}), 6=origen
\NewDocumentCommand{\PreguntaTestSingle}{m m m m m m}{%
  \par\medskip
  \noindent\textbf{#1.}~#2\par\smallskip

    % Imagen (si no es {})
    \IfBlankTF{#5}{}{%
      \begin{center}
        \includegraphics[width=0.75\linewidth]{#5}
      \end{center}
    }

    \begin{Form}
      #3
    \end{Form}

    \par\smallskip
    \switchocg{sol-#1}{\fbox{Mostrar / ocultar solución}}%
    \hfill{\footnotesize #6}\par\smallskip

    \begin{ocg}{Solución}{sol-#1}{0}
      \noindent #4
    \end{ocg}
  \par\medskip
}

% Pregunta multi-choice (checkboxes)
\NewDocumentCommand{\PreguntaTestMulti}{m m m m m m}{%
  \par\medskip
  \noindent\textbf{#1.}~#2\par\smallskip

    % Imagen (si no es {})
    \IfBlankTF{#5}{}{%
      \begin{center}
        \includegraphics[width=0.75\linewidth]{#5}
      \end{center}
    }
    \begin{Form}
      #3
    \end{Form}

    \par\smallskip
    \switchocg{sol-#1}{\fbox{Mostrar / ocultar solución}}%
    \hfill{\footnotesize #6}\par\smallskip

    \begin{ocg}{Solución}{sol-#1}{0}
      \noindent #4
    \end{ocg}
  \par\medskip
}

%% ------------------- Listas----------------------%%
    % --- Lista alfabética ---
    \newenvironment{la}
      {\begin{list}{\alph{alphaitem})}{%
          \usecounter{alphaitem}%
          \setlength{\leftmargin}{1cm}%
          \setlength{\parskip}{3pt}% 
          \setlength{\itemsep}{0pt}% ← espacio entre ítems
          \setlength{\parsep}{0pt}%  ← párrafos dentro de un ítem
      }}
      {\end{list}}
      
    % --- Lista numérica ---
    \newenvironment{lnum}
      {\begin{list}{\arabic{numitem}.}{%
          \usecounter{numitem}%
          \setlength{\leftmargin}{1cm}%
          \setlength{\parskip}{3pt}% 
          \setlength{\itemsep}{0pt}% ← espacio entre ítems
          \setlength{\parsep}{0pt}%  ← párrafos dentro de un ítem
      }}
      {\end{list}}

%% ------------------- Preguntas Test ----------------------%%
      \usepackage{xparse} % si ya lo tienes, no lo repitas

    \NewDocumentEnvironment{test}{m m o}{%
      \begin{tcolorbox}[
        colback=white,
        colframe=black!15,
        boxrule=0.6pt,
        arc=2mm,
        left=2mm,right=2mm,top=2mm,bottom=2mm
      ]
      \centering
      \includegraphics[width=\linewidth]{#1}
      \end{tcolorbox}
    
      \vspace{2mm}
    
      % Caja SOLO para texto (SÍ partible y mantiene márgenes al partir)
      \begin{tcolorbox}[
        enhanced,
        breakable,
        colback=black!3,
        colframe=black!25,
        boxrule=0.6pt,
        arc=2mm,
        left=2mm,right=2mm,top=1.5mm,bottom=1.5mm
      ]
      \textbf{Respuesta:} \textbf{#2}%
      \IfValueT{#3}{\par\smallskip \textbf{Justificación:} #3}
      \end{tcolorbox}
    }{}
      
% ------------------- Imágenes----------------------%
    \graphicspath{{Img/}}% Rutas por defecto 
    % --- Imagen adaptativa ---
    % Registros de dimensión definidos una sola vez
    \newdimen\imgcap
    \newdimen\imgmin
    \newdimen\imgmax
    \newdimen\imgremain
    \newdimen\imgh
    \newdimen\imgneed
    
    \newcommand{\imagenfit}[3]{%
      \addvspace{10pt}% separación antes
      \par\begingroup
        % Parámetros de composición (asignaciones, no \newdimen)
        \imgcap=1\baselineskip            % alto estimado de título+fuente
        \imgmin=0.15\textheight
        \imgmax=0.15\textheight
        \imgremain=\pagegoal
        \ifdim\pagegoal=\maxdimen \imgremain=0pt\fi
        \advance\imgremain by -\pagetotal
        \advance\imgremain by -\imgcap
        % Decisión de altura destino
        \ifdim\imgremain<\imgmin
          \newpage
          \imgh=0.20\textheight
        \else
          \imgh=\imgremain
          \ifdim\imgh>\imgmax \imgh=\imgmax \fi
        \fi
        % Reservar espacio para evitar cortes del bloque completo
        \imgneed=\imgh \advance\imgneed by \imgcap
        \Needspace{\imgneed}
        % Cuerpo
        \noindent\begin{minipage}{\textwidth}
          \centering
          \captionof{figure}{\textemdash\ #2}\par\vspace{0.3em}% título arriba
          \includegraphics[width=\linewidth,height=\imgh,keepaspectratio]{\detokenize{#1}}\par
          \vspace{0.3em}\small\textit{Fuente: }#3% fuente abajo
        \end{minipage}%
      \endgroup
      \par\addvspace{10pt}%
    }

%------------ Bloque de RECUADRO ESPECIAL -------------%
    \newenvironment{specialblock}[1]{%
      \begin{quote} % crea sangría a izquierda y derecha
      \small % texto más pequeño
      \noindent\textbf{#1}\\[-0.3em] % título en negrita
      \rule{\linewidth}{0.4pt}\\[0.6em]}{
      \end{quote}}
      
%-------------------------- Portada -----------------------%
\newcommand{\oldtitle}[1]{\def\OldTitle{#1}}
\newcommand{\changerationale}[1]{\def\ChangeWhy{#1}}

\newcommand{\changebox}{%
\begin{tcolorbox}[title={Historial de cambios del tema}]
\textbf{Título anterior:} \OldTitle\par
\textbf{Motivo del cambio:} \ChangeWhy
\end{tcolorbox}
}

%-------------------------- Author Font -----------------------%
    \newcommand{\authorfont}[1]{{\fontfamily{EBGaramond-TLF}\selectfont #1}}

%-------------------------- Anexos -----------------------%
    \makeatletter
    \newcounter{anexo}
    \renewcommand{\theanexo}{\Roman{anexo}}
    \newcommand{\anexo}[1]{%
      \clearpage
      \phantomsection
      \refstepcounter{anexo}%
      \noindent\textbf{Anexo \theanexo.- }#1\par\medskip
      \addcontentsline{stc}{section}{Anexo \theanexo.- #1}%
    }
    \makeatother

    %-------------------------- Lista de figuras (personal) -----------------------%
\makeatletter
\newcommand{\MyListOfFigures}{%
  \clearpage
  \phantomsection
  {\centering\bfseries\Large Índice de figuras\par}%
  \medskip
  % Entrada en nuestro índice de contenido (stc)
  \addcontentsline{stc}{section}{Índice de figuras}%
  % Recuperación del listado (sin cabecera propia)
  \begingroup
    \parindent=0pt\parskip=2pt
    \@starttoc{lof}%
  \endgroup
}
\makeatother

%-------------------------- Bibliografía -----------------------%
\makeatletter
% Contador para las referencias
\newcounter{myref}
% Cabecera de la bibliografía, entra en el índice 'stc'
\newcommand{\MyBibliography}{%
  \clearpage
  \phantomsection
  {\centering\bfseries\Large Bibliografía y referencias\par}%
  \medskip
  \addcontentsline{stc}{section}{Bibliografía y referencias}%
  \setcounter{myref}{0}%
}
\newcommand{\MyBibItem}[4]{%
  \refstepcounter{myref}%
  \noindent[\themyref]\ #1.\ \emph{#2}. #3. #4\par
}
\makeatother

%--------- Establecer cuatro niveles de índice --------%
    \setcounter{tocdepth}{4}   
    \setcounter{secnumdepth}{4}% numera hasta \paragraph

%%%%%%%%%%%%%%%%%%%%%%%%%%%%%%%%

\makeatletter

    %--------- Interlineado Global --------%
    \edef\GlobalBaseLineStretch{\baselinestretch}
    
    %--------- Espaciado de seccinones --------%
    \renewcommand\section{\@startsection{section}
      {1}{\z@}
      {2.5ex \@plus 1ex \@minus .2ex} % espacio antes
      {1.5ex \@plus .2ex}             % espacio después
      {\normalfont\Large\bfseries}    % formato del título
    }
    \renewcommand\subsection{\@startsection{subsection}{2}{\z@}
      {3ex \@plus 0.8ex \@minus .2ex}
      {1.5ex \@plus .2ex}
      {\normalfont\large\bfseries}
    }
    \renewcommand\subsubsection{\@startsection{subsubsection}{3}{\z@}
      {3ex \@plus 0.5ex \@minus .2ex}
      {1ex \@plus .2ex}
      {\normalfont\normalsize\bfseries}
    }
    \renewcommand\paragraph{\@startsection{paragraph}{4}{\z@}%
    {-2ex \@plus -0.5ex \@minus -.2ex}% espacio antes
    {0.8ex \@plus .2ex}% espacio después
    {\normalfont\normalsize\itshape}}% ← cursiva en lugar de negrita

    %--------- Ecuaciones --------%   
    % Variables globales para almacenar títulos y notas
    \gdef\leq@current@section{}%
    \gdef\leq@current@subsection{}%
    \gdef\leq@last@printed@section{}%
    \gdef\leq@last@printed@subsection{}%
    
    % Comando para añadir nota (invisible en el cuerpo)
    \newcommand{\eqnote}[1]{%
      \addcontentsline{leq}{leqnote}{#1}%
    }
    
    % Comando para ecuaciones
    \newcommand{\eqblock}[2]{%
      % Verificar si necesitamos imprimir el título de sección
      \ifx\leq@current@section\leq@last@printed@section\else
        \addcontentsline{leq}{leqsection}{\leq@current@section}%
        \global\let\leq@last@printed@section\leq@current@section
        \gdef\leq@last@printed@subsection{}% Reset subsección al cambiar sección
      \fi
      % Verificar si necesitamos imprimir el título de subsección
      \ifx\leq@current@subsection\leq@last@printed@subsection\else
        \ifx\leq@current@subsection\@empty\else
          \addcontentsline{leq}{leqsubsection}{\leq@current@subsection}%
          \global\let\leq@last@printed@subsection\leq@current@subsection
        \fi
      \fi
      % Ecuación
      \refstepcounter{eqnum}%
      \begin{equation*}
        #1\tag{E.\theeqnum}
      \end{equation*}%
      \addcontentsline{leq}{leqt}{[E.\theeqnum] -- #2}%
      \addcontentsline{leq}{leqm}{#1}%
    }
    
    %%% ----- Índice de ecuaciones ----------%%%
    % Tamaño para []
    \newcommand{\leqIndexFont}{\normalsize}
    \newcommand{\leqTitleFont}{\tiny}
    \newcommand{\leqNoteFont}{\tiny}
    % Cabecera de sección 
    \newcommand*\l@leqsection[2]{%
      \addvspace{25pt}% Mayor separación antes de nueva sección
      \noindent\leqIndexFont
      \makebox[\linewidth]{\rule{.38\linewidth}{0.4pt}\ \ \textbf{  \MakeUppercase{#1}}\ \ \rule{.38\linewidth}{0.4pt}}%
      \par\vspace{-10pt}
    }
    % Cabecera de subsección 
    \newcommand*\l@leqsubsection[2]{%
      \par\addvspace{15pt}%
      \noindent\leqIndexFont #1
      \par\vspace{-7pt}
    }
    % Título de ecuación 
    \newcommand*\l@leqt[2]{%
    \par\addvspace{10pt}%
      \par\noindent\leqTitleFont #1\par
    }
    % Miniatura de ecuación
    \newcommand*\l@leqm[2]{%
      \vspace{0.3\baselineskip}%
      \par\scriptsize\noindent\hspace*{1.5em}\(#1\)\par%
    }
    % Nota de ecuación 
    \newcommand*\l@leqnote[2]{%
      \vspace{0.2\baselineskip}%
      \par\noindent\leqNoteFont\hspace*{1.5em}\textbf{Nota.--} #1\par\vspace{0.5\baselineskip}%
    }
    % Generación del índice
    \newcommand{\mylistofequations}{%
      \clearpage
      \phantomsection
      % Título visual de la lista (ajústalo a tu gusto):
      {\centering\bfseries\Large Índice de ecuaciones\par}
      % Entrada en nuestro índice 'stc':
      \addcontentsline{stc}{section}{Índice de ecuaciones}%
      % Llamada a tu lista real (usa el nombre exacto de tu macro):
      \@starttoc{leq}%
    }
    
    %--------- Captura de títulos de sección --------%
    \let\leq@original@section\section
    \let\leq@original@subsection\subsection
    % Flag para saber si estamos en una sección asterisco
    \newif\ifLeqStarSection
    \LeqStarSectionfalse
    
    % Redefinir \section CON CONTROL DE ASTERISCO
    \renewcommand{\section}{%
      \@ifstar{\leq@section@star}{\@dblarg\leq@section@nostar}%
    }
    \def\leq@section@star#1{%
      \global\LeqStarSectiontrue
      \gdef\leq@current@section{#1}%
      \gdef\leq@current@subsection{}%
      \leq@original@section*{#1}%
      \global\LeqStarSectionfalse
    }
    \def\leq@section@nostar[#1]#2{%
      \global\LeqStarSectionfalse
      \gdef\leq@current@section{#2}%
      \gdef\leq@current@subsection{}%
      \leq@original@section[#1]{#2}%
    }
    % Redefinir \subsection
    \renewcommand{\subsection}{%
      \@ifstar{\leq@subsection@star}{\@dblarg\leq@subsection@nostar}%
    }
    \def\leq@subsection@star#1{%
      \global\LeqStarSectiontrue
      \gdef\leq@current@subsection{#1}%
      \leq@original@subsection*{#1}%
      \global\LeqStarSectionfalse
    }
    \def\leq@subsection@nostar[#1]#2{%
      \global\LeqStarSectionfalse
      \gdef\leq@current@subsection{#2}%
      \leq@original@subsection[#1]{#2}%
    }

    % ----- Restauración de valores Globales de estilo ----------%
    % Flags
    \newif\ifInFootnote
    \InFootnotefalse
    \pretocmd{\@footnotetext}{\InFootnotetrue}{}{}
    \apptocmd{\@footnotetext}{\InFootnotefalse}{}{}
    % Profundidad de listas (soporta anidadas)
    \newcount\MyListDepth \MyListDepth=0
    \AtBeginEnvironment{la}{\advance\MyListDepth by 1}
    \AfterEndEnvironment{la}{\advance\MyListDepth by -1}
    \AtBeginEnvironment{lnum}{\advance\MyListDepth by 1}
    \AfterEndEnvironment{lnum}{\advance\MyListDepth by -1}
    \AtBeginEnvironment{itemize}{\advance\MyListDepth by 1}
    \AfterEndEnvironment{itemize}{\advance\MyListDepth by -1}
    \AtBeginEnvironment{enumerate}{\advance\MyListDepth by 1}
    \AfterEndEnvironment{enumerate}{\advance\MyListDepth by -1}
    % Restauración diferida: hazlo al INICIO del próximo párrafo real
    \newif\if@pendingRestore
    \@pendingRestorefalse
    \newcommand{\RequestRestore}{\global\@pendingRestoretrue}
    \everypar{%
      \if@pendingRestore
        \global\@pendingRestorefalse
        \ifInFootnote\else
          \ifnum\MyListDepth=0
            \setlength{\parskip}{\GlobalParskip}%
            \setstretch{\GlobalBaseLineStretch}%
          \fi
        \fi
      \fi
    }
    % Al final de bloques:
    \AfterEndEnvironment{equation}{\RequestRestore}
    \AfterEndEnvironment{equation*}{\RequestRestore}
    \AfterEndEnvironment{la}{\RequestRestore}
    \AfterEndEnvironment{lnum}{\RequestRestore}
    % Al final de macros:
    \apptocmd{\eqblock}{\RequestRestore}{}{}
    \apptocmd{\imagenfit}{\RequestRestore}{}{}
    
\makeatother

% ===================== ToC propio con 4 niveles (único bloque) =====================
\makeatletter

% --- Escritores: añadimos una línea al .stc en cada cabecera numerada ---
% Guardamos las versiones actuales para llamar al comportamiento normal
\let\MyTOC@orig@section\section
\let\MyTOC@orig@subsection\subsection
\let\MyTOC@orig@subsubsection\subsubsection
\let\MyTOC@orig@paragraph\paragraph

% SECTION
\renewcommand{\section}{%
  \@ifstar{\MyTOC@section@star}{\@dblarg\MyTOC@section@nostar}%
}
\def\MyTOC@section@star#1{%
  \MyTOC@orig@section*{#1}% sin entrada al .stc
}
\def\MyTOC@section@nostar[#1]#2{%
  \MyTOC@orig@section[{#1}]{#2}% comportamiento normal (escribe al .toc)
  % Escribimos también nuestra línea al .stc (texto con \numberline)
  \addcontentsline{stc}{section}{\protect\numberline{\thesection}#2}%
}

% SUBSECTION
\renewcommand{\subsection}{%
  \@ifstar{\MyTOC@subsection@star}{\@dblarg\MyTOC@subsection@nostar}%
}
\def\MyTOC@subsection@star#1{%
  \MyTOC@orig@subsection*{#1}%
}
\def\MyTOC@subsection@nostar[#1]#2{%
  \MyTOC@orig@subsection[{#1}]{#2}%
  \addcontentsline{stc}{subsection}{\protect\numberline{\thesubsection}#2}%
}

% SUBSUBSECTION
\renewcommand{\subsubsection}{%
  \@ifstar{\MyTOC@subsubsection@star}{\@dblarg\MyTOC@subsubsection@nostar}%
}
\def\MyTOC@subsubsection@star#1{%
  \MyTOC@orig@subsubsection*{#1}%
}
\def\MyTOC@subsubsection@nostar[#1]#2{%
  \MyTOC@orig@subsubsection[{#1}]{#2}%
  \addcontentsline{stc}{subsubsection}{\protect\numberline{\thesubsubsection}#2}%
}

% PARAGRAPH
\renewcommand{\paragraph}{%
  \@ifstar{\MyTOC@paragraph@star}{\@dblarg\MyTOC@paragraph@nostar}%
}
\def\MyTOC@paragraph@star#1{%
  \MyTOC@orig@paragraph*{#1}%
}
\def\MyTOC@paragraph@nostar[#1]#2{%
  \MyTOC@orig@paragraph[{#1}]{#2}%
  % Si no numeras los \paragraph, cambia \theparagraph por texto plano sin \numberline
  \addcontentsline{stc}{paragraph}{\protect\numberline{\theparagraph}#2}%
}

% --- Impresor del índice propio (lee el .stc) ---
\newcommand{\MyTOC}{%
  \par\bigskip
  {\centering\bfseries\Large Índice de contenido\par}%
  \medskip
  \begingroup
    \parindent=0pt\parskip=2pt
    \@starttoc{stc}%
  \endgroup
}

\makeatother
% ===================== fin ToC propio =====================


%%%%%%%%%%%%%%


% --- Datos del tema ---
\title{Tema 3.A.12 -- Teoría de los costes.\\
\large Análisis de dualidad en el ámbito de la empresa. Aplicaciones empíricas.}
\author{Marcelo Ortega}
\date{Febrero 2026}
\newcommand{\TemaCode}{3A12}

\oldtitle{Tema 3.A.10. Teorías de los costes. Análisis de dualidad.}
\changerationale{
    Se añade una mención explícita al trabajo empírico en la materia que mediante diversas técnicas de estimación se utiliza para construir isocuantas y estudiar la eficiencia empresarial.
}

\usepackage{soul}\sethlcolor{yellow}% resaltado amarillo: \hl{texto} (Panel TCEE, ⌃H)
\begin{document}


\maketitle
\changebox

\newpage

% --- Página de resumen y modificaciones ---
\puntosclave{ %% Escribir puntos clave Apuntes CECO
    }
\vspace{1cm}

\modificaciones{
    
    }

\newpage
\MyTOC
\newpage

\mylistofequations
\clearpage

\MyListOfFigures
\clearpage


\pagenumbering{arabic}
\setcounter{page}{1}


\section{Introducción}\label{intro}


\subsection{Contextualización}\label{context}


\subsubsection{Relevancia}\label{importance}


\subsubsection{Historia}\label{history}



\subsubsection{Problemática}\label{problem} 




\subsection{Estructura de la exposición}\label{structure}







\clearpage
\section{Función de costes: Análisis de dualidad}
\puntosclave{
	El problema dual consiste en minimizar el coste condicionado a un nivel de producción.
    
	Del problema dual se obtiene la demanda condicionada de factores que refleja la variación en la cantidad demandada de los factores ante alteraciones en sus precios, manteniendo el nivel de producción constante.
    
	Evaluando la demanda condicionada de factores en la función objetivo del problema dual se obtiene la función de costes. La función de coste cumple una serie de propiedades: continua, homogenea de grado uno y cóncava en los precios de los inputs, no decreciente respecto a la producción y cumple el Lema de Hotelling
    
	La propiedad de integrabilidad permite recuperar la tecnología de producción de la empresa a partir de la función de costes
    }

La eficiencia técnica no es suficiente para la empresa. Sobre la base de criterios técnicos, podemos descartar procesos o combinaciones productivas técnicamente ineficientes, pero hay muchas formas técnicamente eficientes de utilizar los factores productivos, recogidas en la función de producción. El análisis de los costes introduce la eficiencia económica para seleccionar la forma más barata de acceder a un determinado nivel de producción.


A partir de la teoría de la producción, es decir, del análisis técnico del comportamiento del producto, podemos obtener las isocuantas como fronteras del conjunto de combinaciones de factores que son técnicamente factibles [ver tema 3.A.11]. 

Para ello necesitaremos información económica: los precios de los factores productivos. Dentro de un conjunto de combinaciones técnicamente eficientes (puntos de la isocuanta) resultarán más atractivos para el productor aquéllas combinaciones más intensivas en el factor comparativamente más barato, por lo que tendrá implicaciones esenciales respecto a la demanda de factores productivos por parte de la empresa: resolverá el problema de elegir una combinación de factores que permita alcanzar un nivel de producción determinado tal que su coste sea el más bajo posible.

De este modo, siguiendo a JULIO SEGURA (1996), el problema de optimización del productor se puede plantear de dos maneras:
\begin{la}
    \item Maximización de la producción sujeto a un nivel de costes máximo: $\max_{\{Q(\bar C)\}}$; y,
    \item Minimización del coste sujeto a un nivel de producción mínima: $\min_{\{C(\bar{Q})\}}$.
\end{la}

Por todo ello, los enfoques de maximización y minimización (problemas Duales) parten del mismo conjunto de hipótesis y elementos, pues el problema subyacente es el mismo, sólo cambia la manera de abordarlo. El enfoque dual requiere exactamente la misma especificación del conjunto tecnológico, función de producción y costes que el enfoque primal.

\imagenfit{TeoriaCostes_EsquemaDualidad.png}{Especificación necesaria del problema según el enfoque en el que se plantea. Análisis de dualidad}{Apuntes Víctor Gutiérrez. Tema 3.B.12}

Partimos de un conjunto tecnológico definido en el siguiente marco axiomático o normativo:
\begin{lnum}
    \item \textbf{Axioma 1}. Se trata de un conjunto no vacío;
    \item \textbf{Axioma 2}. Cerrado;
    \item \textbf{Axaioma 3}. Con eliminación gratuita;
    \item \textbf{Axioma 4}. Propiedad de \textit{no free-lunch}; y,
    \item \textbf{Axioma 5}. Posibilidad de inacción
\end{lnum}

Por ahora no definiremos ningún caso concreto de rendimientos a escala, por generalizar el resultado lo máximo posible. En caso de deber ser impuesto, será mencionado exprésamente.

Esto dará lugar a una función de producción, cardinal, de buen comportamiento (lo cual es útil para aplicar calculo diferencial, permitiéndonos plantear el problema del consumidor como un problema de optimización), que tendrá las siguientes características:\footnote{%
    Partimos de la base de que solo incorpora las combinaciones técnicamente eficientes pues consideraremos ahora la Función de Producción con la frontera del Conjunto de Posibilidades de Producción. La eficiencia técnica implica que no existe otro proceso de producción capaz de producir igual cantidad de output utilizando menor cantidad de inputs o mayor cantidad de output utilizando igual cantidad de inputs.
}
\begin{lnum}
    \item No-decreciente en factores productivos: eliminación gratuita;
    \item También será conveniente suponer que existe cierta complementariedad entre los factores productivos, dando lugar a unas isocuantas convexas (estríctamente o no).\footnote{%
        De forma intuitiva, si suponemos dos factores productivos, trabajo y capital, son complementarios en cierto grado si solo podemos producir poco output si uno de los factores es escaso aunque el otro sea abundante. En este sentido, ambos factores productivos son importantes para la producción y la media de dos vectores de producción extremos (uno con capital elevado y trabajo reducido y otro al revés) producirá estrictamente más producto que al menos uno de los dos vectores iniciales extremos (JEHLE y RENY, 2011). ADAM SMITH ya hablaba de complementariedad: los trabajadores son más productivos cuanto más capital haya. Por ejemplo, en una función de producción de tipo Cobb-Douglas existe la idea de complementariedad (la productividad marginal del capital aumenta cuanto aumenta el capital).
    }
\end{lnum}

Dicho esto, hay que tener en cuenta que dicha función será:
\begin{lnum}
    \item Objeto de maximización en el Problema de Maximización de la Producción; y,
    \item Considerada como la restricción en su “Dual”, el Problema de Minimización del Coste.
\end{lnum}

Por otro lado, la empresa hace frente a unos costes, derivados de que los precios de los factores productivos son positivos:\footnote{%
    Vamos a partir siempre de la consideración del coste como coste económico, que tiene en cuenta el coste de oportunidad, en contraposición al coste contable, que no lo tiene en cuenta. Debemos definir qué es el coste de un factor para la empresa. Trabajaremos con el concepto de coste marginal de oportunidad de un factor, que es el valor de la alternativa a la que renuncia la empresa por el uso de una unidad adicional de ese factor. Si la empresa no dispone en este momento de la unidad adicional, tendrá que comprarla (alquilarla) y el coste marginal de oportunidad será el precio de mercado o de alquiler de dicho factor. Si la empresa ya dispone en este momento de la unidad adicional, no tendrá que realizar un desembolso, pero, dado que dicha unidad podía haberse vendido en el mercado, este precio de mercado será el valor de la alternativa que se ha perdido. En resumen, el coste marginal de oportunidad es el precio de mercado.
}
\begin{lnum}
    \item Que serán la restricción del $\max_{\{Q(\bar C)\}}$; y,
    \item Y objeto de minimización en el $\min_{\{C(\bar Q)\}}$
\end{lnum}

NOTA.- Existen varias “acepciones” en la aproximación a la Relación Dual. Por ello, introduciremos primero los tres problemas y luego estableceremos las diferencias entre ellos y las relaciones de dualidad que presentan. Con ello, se podrá escoger la que más convenga para la exposición.

\subsection{Problemas de Maximización}
Cuando estamos maximizando la función de producción, es decir, llevando a cabo el PMaxP, lo que en realidad queremos contestar, dada la naturaleza de la empresa es: ¿cuál será el nivel de producción que maximice nuestros beneficios?

Esta pregunta nos supone un enorme problema, puesto que, como sabemos, para establecer las relaciones de dualidad, las aproximaciones de minimización y maximización deben desarrollarse en el mismo marco normativo y debe asegurarse un trasvase correcto de información entre uno y otro. Ahora bien, el beneficio es un elemento que no está en el $\min_{\{C(\bar Q)\}}$ y por ello deberemos tratar el $\max_{\{Q(\bar C)\}}$ con especial cuidado. 
Hay dos posibles vías de resolución para el Problema de Maximización del Beneficio:\footnote{%
    Nótese que hemos cambiado el nombre para aclarar nuevamente los problemas que supone este elemento en las relaciones de dualidad.
}
\begin{la}
    \item La Maximización del Beneficio en una etapa; y,
    \item La Maximización del Beneficio en dos etapas. 
\end{la}

Ahora bien, estas dos distinciones son relevantes si tenemos en cuenta el marco operativo de la empresa. Es decir, la maximización del beneficio no será igual en un mercado que opera en competencia-perfecta, en un mercado donde opera como oligopolio o en un mercado que opera como proveedor monopolístico, entre otras. De hecho, como más adelante veremos, se puede demostrar que en un mercado que opera en competencia-perfecta (empresa precio-aceptante) las relaciones de dualidad se cumplen en cualquiera de los casos. 

Por otro lado, introduciremos el propuesto por JULIO SEGURA que si es, verdaderamente un problema que presenta relaciones de dualidad directas con el $\min_{\{C(\bar Q)\}}$.


\imagenfit{Dif_ProbMAX_PRODUCCION.png}{Enfoques alternativos del problema de maximización del productor}{Elaboración propia}

\subsubsection{$\max_{\{B^o\}}$: Una etapa}
Vamos a considerar que el objetivo final de la empresa es la maximización de sus beneficios, esto es:

\eqblock{
\begin{aligned}
    &\max_{B^o}\;\: \Pi=\underbrace{I(Q)}_{\text{\scriptsize Ingresos}}-\underbrace{CT(Q)}_{\text{\scriptsize Costes Totales}}\\[1em]
    &\underset{\text{\scriptsize Supuestos}}{\Longrightarrow}\quad \boxed{\max_{0\leq \vec z}\;\;\Pi(\vec z)=p\cdot f(\vec z)-\vec w\cdot \vec z}
\end{aligned}
}{Problema de maximización del productor: Una etapa (maximización de beneficios). Función de producción: Análisis de dualidad}
 
Introducimos, no sólo el concepto de costes, sino también el de ingresos: la estructura de ingresos de la empresa corresponde a la de un mercado perfectamente competitivo de bienes, esto es, consideraremos que vende su output (q) a un precio paramétrico (fijo).\footnote{%
    Debemos definir que es el coste de un factor para la empresa. Trabajaremos con coste marginal de oportunidad de un factor, es decir, el valor de la alternativa a la que renuncia la empresa por el uso de una unidad adicional de ese factor:
\begin{la}
    \item Si la empresa no dispone en este momento de la unidad adicional tendrá que comprarla o alquilarla y el coste marginal de oportunidad será el precio de mercado; y,
    \item Si la empresa ya dispone en este momento de la unidad adicional, no tendrá que realizar un desembolso, pero, dado que dcha unidad podría haberse vendido en el mercado, el valor de la alternativa que se ha perdido será el precio de mercado.
\end{la}
} Partiremos de los siguientes supuestos:
\begin{lnum}
    \item Contexto estático con información perfecta y empresa precio-aceptante en ambos mercados (competencia-perfecta);
    \item Tecnología con sus propiedades [\authorfont{Ver Tema 3.A.11}] representable mediante una función de producción de buen comportamiento monoproducto:
    $$q=f(\vec z)$$
    \item De momento supondremos que la empresa opera en el largo plazo, de modo que puede combinar libremente los factores, que serán variables; y,
    \item Producto homogéneo.
\end{lnum}

El empresario ha de elegir el vector de inputs que resuelve directamente el siguiente problema:\footnote{En resumen, el coste marginal de oportunidad es igual al precio de mercado.
}

$$\max_{0\leq \vec z}\;\;\Pi(\vec z)=p\cdot f(\vec z)-\vec w\cdot \vec z$$
 
\paragraph{Condiciones necesarias (CPO’s)}
Dado que se trata de un problema de optimización condicionada con restricciones de desigualdad utilizaremos el Método de los multiplicadores de Lagrange y veremos si el resultado cumple con las condiciones necesarias (Khun-Tucker) y las condiciones suficientes (Hessiana). Se definen como condiciones necesarias (CPO’s):

\eqblock{
\begin{aligned}
    &\text{CPO's}:\\[1em]
    &\left\{\begin{aligned}
        &\text{Estacionariedad}&&\Longrightarrow\;&&\boxed{\frac{\partial \Pi(\vec z^*)}{\partial z_k}=p\cdot f_k(\vec z^*)-w_k\leq 0}\\[0.5em]
        &\text{No-negatividad}&&\Longrightarrow\;&&\boxed{0\leq z_k^*}\\[0.5em]
        &\text{Holgura complementaria}&&\Longrightarrow\;&&\boxed{z_k^*\cdot [p\cdot f_k(\vec z^*)-w_k]=0}
    \end{aligned}\right\}\Longrightarrow\;\forall k=1,2,\dots, n
\end{aligned}
}{Problema del productor (maximización del beneficio): Condiciones de primer orden. Función de producción, análisis de dualidad}

Las condiciones de primer orden implican:
\begin{lnum}
    \item \textbf{Condición de estacionariedad}. La derivada de la función respecto a cada uno de los outputs será igual al beneficio de este por el valor de producido menos el precio de este en el mercado, siempre menor o igual a cero (en el óptimo, se pretende que esto sea igual a cero para todos);
    \item \textbf{Conición de no-negatividad}. Que los outputs sean positivos o iguales a cero (que haya habido producción o no, pero irreversibilidad es imposible y no existe la producción negativa); y,
    \item \textbf{Condición de holgura complementaria}. Que cada output por la diferencia entre el valor adquirido en ingresos (beneficio por output) y el precio pagado por el sea igual a cero.
\end{lnum}

Por tanto, en el óptimo, la empresa contratará un factor productivo hasta que el valor de su PMg sea igual al ingreso que obtiene gracias a su uso. Además, trabajando con las CPO’s obtendremos que la Relación Marginal de Sustitución Técnica (RMST) se iguala al cociente de precios de estos dos factores productivos en el mercado. Esto implica que, estando en el óptimo, si la empresa desea aumentar marginalmente su producción (intensiva), le será indiferente qué factor productivo contratar, puesto que la PMg de todos ponderadas por el coste es la misma en todos los casos. 

\paragraph{Condiciones suficientes (CSO’s)}
Sin embargo, las condiciones necesarias no son suficientes para garantizar que el resultado sea solución al problema. Para que el resultado del problema de optimización sea un resultado \textbf{factible}, este debe cumplir las condiciones suficientes, que garantizan que nuestro vector es la solución a este problema:

\eqblock{
\begin{aligned}
    H=\left(\frac{\partial^2\Pi(\vec z^*)}{\partial z_k\partial z_l}\right)_{\forall k,l}=\left(p\cdot \frac{\partial^2 I(\vec z^*)}{\partial z_k\partial z_l}\right)_{\forall k,l}=(p\cdot f_{k,l})_{\forall k,l}=p^n\cdot 
    \underbrace{\left(\begin{matrix}
        f_{11}&\cdots &f_{1,n}\\
        \vdots&\ddots&\vdots\\
        f_{n1}&\cdots&f_{nn}
    \end{matrix}\right)}_{F\equiv \text{Función de Producción}}
\end{aligned}
}{Problema del productor (maximización del beneficio): Condiciones Sufucientes. Función de producción, análisis de dualidad}

Es decir, que la matriz Hessiana sea semidifinida negativa. Como los precios son estríctamente positivos, es necesario que $F$ sea semidefinida negativa. En otras palabras, es necesario que la función de producción sea estríctamente cóncava, condición que hemos impuesto de partida.\footnote{%
    Como sobre este problema de optimización no se han impuesto restricciones, si la función de producción fuese convexa, el valor máximo de la función de beneficios se obtendrá cuando tienda la función de producción a infinito, en cuyo caso obtendrá también un beneficio infinito.
}  

Veamos una explicación más detallada sobre esto: las condiciones suficientes requieren estricta concavidad (rendimientos decrecientes a escala) de la función de producción, por lo menos:

\eqblock{
\begin{aligned}
    &F\to\text{CSO}(\checkmark)  \iff&&\underbrace{\text{o,}\quad \frac{\partial CMg}{\partial q}>0;\quad \text{o,}
    \begin{cases}
        \frac{\partial {PMg}_L}{\partial L}<0\\[0.5em]
        \frac{\partial {PMg}_K}{\partial K}<0
    \end{cases}
    }_{\text{\scriptsize Alternativo. Garantiza concavidad función de producción $\sim$ RDecrecientes a Escala}}
\end{aligned}
}{Problema del productor (maximización del beneficio): Propiedades de la Función de Producción y relación con la CSO's. Función de producción, análisis de dualidad}

Para garantizar estas condiciones, los requisitos son más estrictos que para el problema $\min_{\{C(\bar Q)\}}$. En efecto, en diversas situaciones una función de costes no basta para asegurar una solución finita a la maximización de beneficios:
\begin{lnum}
    \item Si el CMe es siempre inferior al precio, los beneficios crecen continuamente con el output;
    \item Si los CMe son decrecientes, y por tanto el CMg está por debajo del medio, igualar precio a coste marginal implica necesariamente pérdidas. Pero esta presencia de economías de escala es una poderosa fuerza hacia posiciones de poder de mercado (en el límite tendríamos el monopolio natural) y por lo tanto dificilmente encajaría en un contexto competitivo;
    \item Si el CMg es decreciente, igualar precio a coste marginal implicaría minimizar beneficios, para una empresa precio-aceptante;
    \item Los rendimientos constantes a escala plantean tres posibilidades:
    \begin{la}
        \item El precio es mayor al CMe (igual que el caso 1), con beneficios positivos, que crecen indefinidamente con el nivel de producción. A nivel de empresa plantean la pregunta de por qué no entrarán más empresas, pero si todas tienen acceso a la misma tecnología, la entrada sólo trasladaría el problema a nivel industria. Realmente la cuestión radica en si los rendimientos pueden mantenerse constantes ilimitadamente, pese a la posibilidad de “replicar” las empresas los procesos seguidos;
        \item Si el precio es inferior al CMe, los beneficios se maximizan con un nivel de output nulo; o,
        \item Si el precio es igual al CMe, con beneficios nulos para cualquier nivel de output: a nivel de industria la demanda de mercado puede resolver la indeterminación, pero quedaría pendiente como la producción se distribuye entre un número de empresas a priori también indeterminado; el tamaño del mercado delimita el resultado de multiplicar el output de cada empresa por el número de empresas, pero no puede decidir cada uno de los factores de esta multiplicación.
    \end{la}
\end{lnum}

En todo caso, en la práctica, un criterio adicional a considerar es la no-negatividad de los  beneficios, ni siquiera a C/p para la cobertura de los costes variables. Por eso hay que añadir a las condiciones anteriores la restricción adicional de que los beneficios (económicos) sean mayores o iguales a cero.  

\paragraph{Sistema Completo de Ecuaciones de Demanda de Inputs No Condicionadas (SCEDINC)}
A partir de este problema de optimización, siempre y cuando tenga solución, podemos obtener el SCDEINC es el conjunto de funciones que resultan del problema de maximización de beneficios de una etapa, donde la empresa elige simultáneamente qué producir y cuánto usar de cada factor. El nivel de output se determina por tanto exógenamente al maximizar los beneficios. Por tanto, es el conjunto de demandas de factores derivadas del problema de maximización y tiene la siguiente forma:

\eqblock{
\begin{aligned}
    \vec Z^*(\vec w, p)=
    \left(\begin{matrix}
        z_1(w_1, w_2, \dots, w_n, p)\\
        \vdots\\
        z_{n}(w_1,w_2, \dots, w_n, p)
    \end{matrix}\right)
\end{aligned}
}{Sistema Completo de Ecuaciones de Demanda de Inputs No Condicionadas (SCEDINC). Función de Producción: Análisis de dualidad}
 
Cada una expresa la cantidad óptima del insumo $i$ como función del precio de los factores y del precio del producto. Posee las siguientes propiedades:
\begin{lnum}
    \item Las funciones de demanda son continuas y homogeneas de grado cero en precios y costes (p, w);\footnote{%
        Si los precios se multiplican por cualquier escalar, la demanda óptima de insumos no varía. Solo las relaciones relativas de precios importan, no sus niveles absolutos.
    } y,
    \item El SCEDINC es, con cierta generalidad, diferenciable tanto en coste de insumos como en precios de producto. 
\end{lnum}

Quedaría representado a través del siguiente gráfico:

\imagenfit{SCEDINC_IB_FP_CSO.png}{Representación de la solución al problema de maximización del beneficio en caso de que se cumplan las CPO's y CSO's. SCEDINC}{Apuntes Víctor Gutiérrez. Tema 3.A.11}

Donde: (i) $IB^{0}$ representa la isobeneficio de nivel ($0$); (ii) $E$ representa el óptimo demandado de inputs; y, ($3$) $f(\vec z)$ representa la función de producción estríctamente cóncava ($f''(\vec z)<0$).\footnote{%
    La pendiente de la recta de isobeneficios es $w/p$ porque así se deriva de la ecuación que hemos maximizado. Es decir, supongamos que a la ecuación a maximizar le asignamos un valor constante $B^{0,1,2,\dots,n}$. Cada recta corresponderá a un nivel de beneficios, entonces, reformulando la ecuación $q$ en función de $z$ (tal y como se aprecia en el gráfico), obtendremos que el cociente multiplicando a $z$ es, precisamente, la ratio de coste de insumos entre precios.
}

\paragraph{La Función de Oferta (no condicionada) de output}
La función de oferta no condicionada de output es la ecuación que pone la cantidad de output ofertado en función del coste de los insumos y el precio del producto, es decir, con ella podremos obtener fácilmente el output en el óptimo.  

\eqblock{
\begin{aligned}
    &\text{Función de oferta no condicionada de output:}&&\boxed{q^*(\vec w, p)=f(\vec z^*)}
\end{aligned}
}{Función de Oferta (no condicionada de output). Función de Producción: Análisis de dualidad}
 
Esta es continua y homogenea de grado cero tanto en precios como en costes ($p, w$).

\paragraph{Función Valor de Beneficios}
Se trata de la función de valor máximo que obtenemos al sustituir las soluciones del problema anterior (es decir, al sustituir las demandas no condicionadas de los factores y la función de oferta no condicionada de output en la función objetivo de beneficios):

\eqblock{
\begin{aligned}
    &\text{Función Valor}&&\Pi^*=p\cdot f(\vec z)-\vec w\cdot \vec z^*=\Pi^*(\vec w, p)
\end{aligned}
}{Función Valor de beneficios. Función de Producción: Análisis de dualidad}

La función de beneficios posee las siguientes propiedades:
\begin{lnum}
    \item Es continua y homogénea de grado cero tanto en precios como en costes ($p,w$);
    \item Es estríctamente creciente en precios;
    \item Es no-creciente en costes; y,
    \item Es convexa tanto en precios como en costes ($p,w$); y,
    \item Cumple el Lema de Hotelling:
    $$\frac{\partial\Pi(p,w)}{\partial p}=q(p,w)\qquad \text{y}\qquad -\frac{\partial \Pi(p,w)}{\partial w_i}=z_i(p,w)$$
    La derivada del beneficio máximo respecto al precio del producto da la oferta óptima de consumo; la derivada del beneficio máximo respecto al coste de un factor (con signo negativo) da la demanda óptima de ese factor. Es decir, las funciones de oferta de output y demanda de inputs pueden obtenerse directamente a partir de la función de beneficio máximo.\footnote{%
        Como la función de beneficio máximo ya incorpora el comportamiento óptimo de la empresa;un pequeño cambio en precios o en coste de factores, solo afecta al beneficio a través de cómo cambio el ingreso o el costo en el punto óptimo, no por variaciones en las cantidades (que ya están optimizadas).
    }
\end{lnum}

Estas dos últimas propiedades llevan a reconocer ciertas propiedades de la matriz de sustitución de la función de beneficios, la matriz Hessiana de la función de beneficios respecto de todos los precios y costes:

\eqblock{
\begin{aligned}
    &\text{Matriz de sustitución:}&&\Xi\equiv
    \left(\begin{matrix}
        \partial^2\Pi/\partial p^2&\partial^2\Pi/\partial p\partial w_1&\cdots &\partial^2\Pi/\partial p\partial w_1\\
        \partial^2\Pi/\partial w_1\partial p&\partial^2\Pi/\partial w_1^2&\cdots &\partial^2\Pi/\partial w_1\partial w_n\\
        \vdots &\vdots &\ddots &\vdots\\
        \partial^2\Pi/w_n\partial p&\partial^2\Pi/\partial w_n\partial w_1&\cdots &\partial^2\Pi/\partial w_n^2\\
    \end{matrix}\right)\\[1em]
    &\underset{\substack{\text{\scriptsize Por el Lema}\\\text{\scriptsize de HOTELLING}}}{\Longrightarrow} && \Xi=
    \left(\begin{matrix}
        \partial q^*/\partial p&\partial q^*/\partial w_1&\cdots &\partial q^*/\partial w_n\\
        -(\partial z_1^*/\partial p)&-(\partial z_1^*/\partial w_1)&\cdots &-\partial (z_n^*/\partial w_n)\\
        \vdots &\vdots &\ddots &\vdots\\
        -(\partial z_n^*/\partial p)&-(\partial z_n^*/\partial w_1)&\cdots &-(\partial z_n^*/\partial w_n)\\
    \end{matrix}\right)
\end{aligned}
}{Matriz de sustictución de la Función de Beneficios: Matriz Hessiana de $\Pi$ respesto $(p,w_i)_{\forall i}$. Función de Producción: Análisis de dualidad}
 
Por tanto, la matriz de sustitución será:
\begin{lnum}
    \item Simétrica, de modo que el orden de derivación no importa, lo que implica que:
    \begin{la}
        \item Las demandas de inputs son simétricas en el precio de los factores:
        $$\partial^2\Pi/\partial w_k\partial w_l=\partial^2\Pi/\partial w_l\partial w_k\iff \partial z_k^*/\partial w_l=\partial z_l^*/\partial w_k$$
        \item El comportamiento de la oferta de output es simétrico (pero opuesto) al de cada una de las demandas de input:
        $$\partial^2\Pi/\partial w_k\partial p=\partial^2\Pi/\partial p\partial w_k\iff-(\partial z_k^*/\partial p)=\partial q^*/\partial w_k$$
    \end{la}
    \item Semidefinida negativa (por la convexidad de la función de beneficios), por lo que los elementos de su diagonal principal son no-negativos, esto es:
    \begin{la}
        \item La curva de oferta de output es no-decreciente en precios:
        $$\partial q^*/\partial p\ge0$$
        \item Las curvas de demanda de inputs son no-crecientes con respecto a su precio:
        $$\partial z_k^*/\partial w_k\leq0$$
    \end{la}
\end{lnum}
 
\subsubsection{$\max_{\{\Pi(\bar C)\}}$: Dos etapas}



\subsubsection{$\max_{\{Q(\bar C)\}}$: SEGURA}



\subsection{Problemas de Minimización}
El Problema de Minimización de Costes fue introducido por PAUL SAMUELSON (1947) y analizado con mayor detalle por RONALD SHEPHARD (1953). 
Conociendo el precio de los factores productivos, podemos encontrar la combinación de factores productivos (eficiencia técnica) capaz de minimizar el coste de producir (eficiencia económica) una determinada cantidad de output para la empresa.

\subsubsection{Supuestos}
Nos encontramos bajo el siguiente marco normativo:
\begin{lnum}
    \item Contexto estático con información perfecta y empresa precio-aceptante en el mercado de inputs (sólo minimiza costes, la estructura de mercado del mercado del output es irrelevante, ya que el precio del bien final no aparece en el problema de optimización);
    \item Tecnología con sus propiedades [Tema 3.A.11] representable mediante una función de producción de buen comportamiento monoproducto:
    $$q=f(\vec z)$$
    \item Supondremos que la empresa opera al largo plazo, de modo que puede combinar libremente los factores; y,
    \item Producto homogéneo.
\end{lnum}

\subsubsection{Desarrollo }
Matemáticamente el problema se puede especificar de la siguiente manera:

\eqblock{
\begin{aligned}
    &\boxed{
    \begin{aligned}
        \min_{\vec z}\quad C(\vec z)=\vec w\cdot \vec z\\
        \text{s.a.}\quad 
        \begin{cases}
            \bar Q\leq f(\vec z)\\
            0\leq \vec z
        \end{cases}
    \end{aligned}
    }\\[2em]
    &\Longrightarrow&&\mathcal L=\vec w\cdot \vec z+\mu(\bar Q-f(\vec z))
\end{aligned}
}{Problema de minimización del gasto. Función de producción: Análisis de dualidad}
 
Dado que se trata de un problema de optimización condicionada con restricciones de desigualdad, nos apoyaremos en el Método de los Multiplicadores de Lagrange para su resolución. El multiplicador nos informa sobre la sensibilidad de la función objetivo (i.e. el coste) ante cambios en la restricción (i.e. en la producción). En concreto, representa lo que varía el coste al variar marginalmente la producción objetivo y su ssentido económico se interpreta como el coste marginal de la producción (${CMg}_Q$).

\subsubsection{Condiciones necesarias (CPOS’s, Khun-Thucker)}
Las condiciones de primer orden o condiciones necesarias para que el vector óptimo de demanda de factores sea la solución a este problema son las siguientes:

\eqblock{
\begin{aligned}
    &\text{CPO's}:\\[1em]
    &\left\{\begin{aligned}
        &\text{Estacionariedad}&&\Longrightarrow\;&&\boxed{\frac{\partial C(\vec z^*)}{\partial z_k}=w_k-\mu^*\cdot f_k(\vec z)\ge 0}\\[0.5em]
        &\text{No-negatividad}&&\Longrightarrow\;&&\boxed{0\leq z_k^*}\\[0.5em]
        &\text{Restricción producción}&&\Longrightarrow\;&&\boxed{f(\vec z^*)=\bar Q}\\[0.5em]
        &\text{Holgura complementaria}&&\Longrightarrow\;&&\boxed{z_k^*\cdot [w_k-\mu^*f_k(\vec z^*)]=0}
    \end{aligned}\right\}\Longrightarrow\;\forall k=1,2,\dots, n
\end{aligned}
}{Minimización del gasto: Condiciones de primer orden. Función de producción, análisis de dualidad}
 
En el óptimo, el multiplicador de Lagrange sería igual al cociente entre el coste del bien y su $PMg$.\footnote{%
    Nótese que se prescinde de la notación extendida, esto es, la PMg del bien (esto es, la derivada de la función de producción respecto al bien) utiliza la notación $f_{k}(z^{*})$
} Es decir, la medida en cómo se traslada la producción a costes depende del coste del factor, ponderado por su productividad. Esto implica que la empresa contrata un factor productivo hasta que su coste marginal es igual a su productividad marginal, cumpliendo la igualdad siguiente:

$$\frac{w_k}{f_k(\vec z^*)}=\mu^*\qquad \forall k=1,2,\dots, n$$
2
 
Y, estableciendo relaciones entre los diferentes factores productivos y el cociente entre sus precios, obtenemos una conclusión ya esperada: (suponiendo únicamente como factores productivos capital y trabajo)

$$\frac{w_K}{f_K(\vec z^*)}=\mu^*=\frac{w_L}{f_L(\vec z^*)}\;\Longrightarrow\;\underbrace{\frac{f_L(\vec z^*)}{f_K(\vec z^*)}}_{|{RMST}_L^K|}=\frac{w_L}{w_K}$$
 
Esto es, trabajando con las CPO’s, obtenemos que la RMST se iguala al cociente de precios de los factores productivos.\footnote{%
    [Tema 3.A.11] La cantidad de un factor productivo al que puede renunciar una empresa si aumenta la cantidad de otro factor productivo en una unidad para mantener el nivel de producción constante.
} Por lo que, en el óptimo, si la empresa desea aumentar marginalmente su producción, le será indiferente contratar más capital que más trabajo porque la PMg ponderada por el CMg es el mismo en ambos casos.

\subsubsection{Condiciones suficientes (CSO’s)}
Para garantizar que el vector de demanda óptimo de insumos es el vector solución a este problema es necesario que se den las CSO’s (condiciones suficientes), consistentes en que los menores del Hessiano orlado correspondiente sean negativos.\footnote{%
    El Hessiano orlado es una matriz que se utiliza para analizar la segunda derivada condicionada de una función objetivo sujeta a una o varias restricciones. Geométricamente, el Hessiano orlado describe la curvatura de la función objetivo restringida a la superficie $f(z) = 0$ [Restricción igual a cero]. En el punto crítico, mide cómo cambia la función objetivo a lo largo de las direcciones tangentes a la restricción. Si el Hessiano orlado es una matriz:
\begin{la}
    \item \textbf{Definida negativa}. Curvatura hacia abajo. La función objetivo decrece en todas las direcciones tangentes por lo que es un máximo restrinigido (local); y,
    \item \textbf{Definida negativa}. Curvatura hacia arriba. La función objetivo crece en todas las direcciones tangentes a la restricción por lo que es un mínimo restringido (local)
\end{la}
} En nuestro caso, como el Lagrangiano está compuesto por una única restricción y la función objetivo es una función compuesto de diferentes inputs utilizados para la producción, el Hessiano orlado vendrá definido por:

\eqblock{
\begin{aligned}
    &\mathcal L=\vec w\cdot \vec z+\mu[\bar Q-f(\vec z)]\\[1em]
    &\underset{\substack{\text{\scriptsize Hessiano}\\\text{\scriptsize orlado}}}{\widehat H}=\left[\begin{matrix}
        0&-\left(\frac{\partial f}{\partial z_1}(z_1^*,z_2^*, \dots, z_n^*,\mu^*)\right)&\cdots & -\left(\frac{\partial f}{\partial z_n}(z_1^*,z_2^*, \dots, z_n^*,\mu^*)\right)\\
        -\left(\frac{\partial f}{\partial z_1}(z_1^*,z_2^*, \dots, z_n^*,\mu^*)\right)&-\left(\frac{\partial^2 \mathcal L}{\partial^2 z_1}(z_1^*,z_2^*, \dots, z_n^*,\mu^*)\right)&\cdots &-\left(\frac{\partial^2 \mathcal L}{\partial z_1\partial z_n}(z_1^*,z_2^*, \dots, z_n^*,\mu^*)\right)\\
        \vdots &\vdots &\ddots&\vdots \\
        -\left(\frac{\partial f}{\partial z_n}(z_1^*,z_2^*, \dots, z_n^*,\mu^*)\right)&-\left(\frac{\partial^2 \mathcal L}{\partial z_n\partial z_1}(z_1^*,z_2^*, \dots, z_n^*,\mu^*)\right)&\cdots &-\left(\frac{\partial^2 \mathcal L}{\partial^2 z_n}(z_1^*,z_2^*, \dots, z_n^*,\mu^*)\right)    
    \end{matrix}\right]\\[1em]
    &\Longrightarrow\qquad\boxed{\det{\widehat H}<0}
\end{aligned}
}{Condiciones suficientes: Propiedades del Hessiano orlado. Función de producción: Análisis de dualidad}

Si el determinante del Hessiano orlado es inferior a cero, entonces podemos garantizar que nos encontramos ante un mínimo.\footnote{%
    En la derivación que proponemos en las Notas finales veremos como la condición es contraria, se exige que sea mayor que cero para un máximo. Esto no es sustancial ya que varía en función de cómo formulemos el Lagrangiano:
$$\mathcal{L}=w\cdot z+\mu(\bar Q-f(z))\qquad \text{o,}\qquad \mathcal L'=w\cdot z-\lambda(f(z)-\bar Q)$$
} Lo más importante es el significado geométrico o matemático: para poder optimizar la producción mediante la resolución del problema de minimización del gasto, la función de producción debe tener alguna región o subespacio cuya estructura a escala sea decreciente o constante, es decir, la función de producción debe tener una región de concavidad (no estricta) para que pueda existir una solución interior óptima.
Supondremos a partir de ahora que la Función de Producción es estríctamente cuasi-cóncava (para facilitar el análisis).
El Hessiano orlado indicaría que el Lagrangiano es convexo, por lo que cualquier desplazamiento factible que mantenga la producción supnodría un aumento del costo. Es decir, el plano isocosto es tangente inferior a la superficie de producción, es tangente en un punto donde la isocuanta es curvada hacia arriba respecto al plano isocosto. 

\textbf{\textcolor{red}{[IMPRESCINDIBLE LA REPRESENTACIÓN GRÁFICA DEL HESSIANO ORLADO PARA PODER ENTENDER EL PROBLEMA]}}

\subsubsection{Desarrollo gráfico}
Gráficamente, para el caso de dos inputs, representaríamos:
\begin{la}
    \item La isocuanta del nivel mínimo de producción que queremos alcanzar; y,
    \item Un mapa de curvas isocostes, que unan los puntos en los que distintas combinaciones de inputs tengan el mismo coste.
\end{la}

\subsubsection{Solución}
El equilibrio se encuentra en el punto de tangencia entre ambas, que supondrá la isocoste más cercana al origen sujeta a la restricción. En este punto, la RMST, es decir, el cociente de productividades marginales de los factores (pendiente de la curva isocuanta) se iguala al cociente del precio de los factores (pendiente de la curva isocoste).

\imagenfit{ISOCUANTAS_PMIN.png}{Punto de tangencia (equilibrio). Curva isocostes y restricción del nivel de producción}{Apuntes Víctor Gutiérrez. Tema 3.B.12}

\paragraph{Sistema Completo de Ecuaciones de Demanda de Inputs Condicionada: SCEDIC}
El SCEDIC describe cómo varían las demandas óptimas de factores productivos (inputs) que una empresa emplea, dado que produce un nivel determinado de output, cuando cambian los precios de los factores:

\eqblock{
\begin{aligned}
    \vec h^*(\vec w, \bar Q)=\left(\begin{matrix}
        h_1(w_1,w_2, \dots, w_n,\bar Q)\\
        \vdots\\
        h_n(w_1,w_2, \dots, w_n,\bar Q)\\
    \end{matrix}\right)
\end{aligned}
}{Problema de Minimización del Gasto: Sistema Completo de Ecuaciones de Inputs Condicionada (SCEDIC). Función de producción: ANálisis de dualidad}

Este Sistema de ecuaciones tiene las siguientes propiedades:
\begin{lnum}
    \item Las funciones de demanda son contínuas y homogéneas de grado cero tanto en costes como en cantidad de output $(p,Q)$; y,
    \item Es, con cierta generalidad, diferenciable tanto en costes como en output
\end{lnum}

\paragraph{Función Valor de Costes}
La Función Valor de costes es la función de valor mínimo que obtenemos al sustituir las soluciones del problema anterior (es decir, al sustituir las demandas condicionadas de los factores en la función objetivo de costes). Es de gran relevancia ya que presupone:
\begin{lnum}
    \item Eficiencia técnica, ya que hemos impuesto que cumpla la función de producción; y,
    \item Eficiencia económica, ya que, al introducir en el análisis el precio de los factores (coste), obtenemos el mínimo al que producir un nivel óptimo. 
\end{lnum}

\eqblock{
\begin{aligned}
    C^*(\vec w, \bar Q)=\vec w\cdot \vec h^*=\vec w\cdot \vec h^*(\vec w, \bar Q)
\end{aligned}
}{Problema de Minimización del Gasto: Función Valor de costes. Función de producción: Análisis de dualidad}

Además, posee las siguientes propiedades:
\begin{lnum}
    \item Es contínua tanto en costes como en cantidades de output;
    \item Es homogenea de grado uno en costes de los factores (w);
    \item Es estríctamente creciente en cantidades de output;
    \item Es no-decreciente, y cóncava, en coste de los factores. Esto se explica por la sustitución entre factores productivos ante variaciones en sus precios (coste). Si el precio de un factor aumenta, el coste aumentaría de forma proporcional si la empresa no modifica la combinación de factores. Sin embargo, cabe esperar que la empresa reaccione y el aumento en los costes aumente menos que proporcionalmente;\footnote{%
        Evidentemente, la tasa en la que se incrementa el coste viene definida por el grado de sustituibilidad de los factores, lo que se recoge en el concepto de elasticidad de sustitución de los factores [ver Tema 3.A.11]
    } y,
    \item Cumple el Lema de Shepard:
    $$h_K^*(\vec w, \bar Q)=\frac{\partial C^*(\vec w, \bar Q)}{\partial w_K}\qquad \forall k$$
\end{lnum}

Estas dos últimas propiedades llevan a reconocer ciertas propiedades de la matriz de sustitución de la función de costes, definida como la matriz Hessiana de la función de costes respecto de todos los precios:

\eqblock{
\begin{aligned}
    \Xi=\left(\begin{matrix}
        \xi_{11}&\cdots &\xi_{1n}\\
        \vdots &\ddots &\vdots \\
        \xi_{n1}&\cdots &\xi_{nn}
    \end{matrix}\right)
    \qquad \text{donde,}\qquad
    \begin{aligned}
        &\xi_{ki}\equiv \frac{\partial C^*(\vec w, \bar Q)}{\partial w_k, \partial w_i}\overbrace{=}^{\substack{\text{\scriptsize Por el Lema}\\\text{\scriptsize de SHEPARD}}}\frac{\partial h_K^*(\vec w, \bar Q)}{\partial w_i}\\[0.5em]
        &\hspace{5em}\text{\scriptsize$\forall k,i=1, 2, \dots, n\;\;\;\;k\neq i$}
    \end{aligned}   
\end{aligned}
}{Problema de Minimización del Gasto: Matriz Hessiana (Condiciones suficientes). Función de Producción: Análisis de dualidad}

La matriz de sustitución es:
\begin{lnum}
    \item Simétrica, de modo que el orden de derivación no importa, es decir, los efectos precio-cruzado para dos factores cualesquiera en las demandas condicionadas son idénticos; y,
    \item Semidefinida negativa (por la concavidad de la función de costes), lo que implica que los elementos de su diagonal principal son no-negativos, es decir, los efectos precio propios en las demandas hicksianas sean no-positivos.\footnote{%
        Esto es así porque todos los elementos de la diagonal principal de una matriz semidefinida negativa deben ser no positivos (nulos o negativos). Y estos elementos de la matriz de sustitución son los efectos precio propios de modo que los efectos precio-propios son negativos, lo que implica que si aumenta el precio de un factor, su demanda no condicionada se reduce o se mantiene constante, pero en ningún caso aumenta.
    }
\end{lnum}

\subsection{Primas vs. Dual}
Una vez que hemos abordado ambos problemas, estamos en condiciones de profundizar en la relación entre ambos. Al igual que hicimos en la teoría del consumo [ver tema 3.A.9], vamos a introducir y sintetizar una serie de resultados que nos permitirán obtener toda Función obtenida en la Teoría de la Producción a partir de cualquier otra.

\imagenfit{Esquema_DUALIDAD_Produccion.png}{Esquema relaciones duales. Un único problema de minimización y diferentes problemas de maximización}{Apuntes Víctor Gutierrez. Tema 3.A.12}
 
\subsubsection{Dualidad entre funciones de costes y la función de producción}
Una aplicación muy útil de la teoría de la dualidad es que ésta nos permite asegurar que la función de costes proviene de la optimización de una función de producción. Es decir, que la función de costes contiene esencialmente la misma información que la función de producción. En particular, si disponemos de la función de costes es posible recuperar la función de producción siguiendo un procedimiento análogo al utilizado en la teoría del consumo [ver tema 3.A.9]:

$$C^*(\vec w, \bar Q)\;\underset{\substack{\text{\scriptsize Invirtiendo, y}\\\bar C\equiv C^*}}{\Longrightarrow}\;q^*=\zeta(\vec w,\bar C)$$
 
La función de la derecha nos informa sobre el máximo output que es posible producir dados los precios de los factores y un nivel de coste dado. A partir de ella puede recuperarse la tecnología original mediante el siguiente problema:\footnote{De la misma forma que hicimos para la integrabilidad de las preferencias. 
}

$$\begin{aligned}
    &\arg \min_{\vec z}\quad \zeta(\vec w, \bar C)\\
    &\text{s.a.}\quad \vec w\cdot \vec z\leq \bar C
\end{aligned}$$

\subsubsection{Dualidad entre la función de beneficio y la de producción}
La tecnología puede representarse mediante el conjunto productivo ($\mathcal Y$), espacio delimitado por la función $f$ y definido como:

$$\underset{(2)}{\\[2em]\mathcal Y=\{(q, \vec z)>0\;;\;q<f(\vec z)\}}$$

Pero, de forma equivalente, la función de Beneficios es la función de soporte de $\mathcal Y$, es decir:

$$\mathcal Y=\{(q,\vec z)>0\;:\;p\cdot q-\vec w\cdot \vec z<\Pi(\vec w, p)\}$$

Gráficamente:

\imagenfit{Rel_DUALES_ConTECNOLOGICO.png}{Conjunto tecnológico y fronteras. Relaciones duales beneficio-gasto}{Apuntes Víctor Gutiérez. Tema 3.A.12}

Como la función de beneficios es homogénea de grado uno en el vector de precios y que éste es estríctamente mayor que cero, entonces:

$$\underset{(3)}{\\[2em]\Pi\left(\frac{1}{p}\cdot p,\frac{1}{p}w\right)=\frac{1}{p}\Pi(p,w)\;\underset{\text{\scriptsize Reales}}{\Longrightarrow}\;\Pi(1, \omega)\quad \text{con}\;\;\;\omega=\frac{w}{p}}$$

Utilizando estas relaciones (2 y 3), dividiendo por $p$ y reordenando tenemos que:

$$\underset{(4)}{\\[1em]\mathcal Y\equiv \{(q, \vec z)>0\;:\;q\leq \min{[\omega\cdot \vec z+\Pi(1, \omega)]\;\;;\;\;\text{s.a.}\;\;\omega \gg0}\}}$$

La expresión a minimizar es convexa (dado que es la suma de dos funciones convexa). Las condiciones necesarias de mínimo nos permiten encontrar sus valores óptimos ($\omega^*$) a partir de las siguientes ecuaciones:

$$\frac{\partial}{\partial w_k}[\omega\cdot \vec z+\Pi(1,\omega)]=z_k+\frac{\partial \Pi(1, \omega)}{\partial w_k}=0;\qquad \forall k=1, 2, \dots, n$$
 
Por el Lema de Hotelling -que la derivada parcial de la función de beneficio respecto a cada uno de los costes de cada factor es igual a la cantidad demandada de éste para el nivel óptimo, cambiado de signo, sustituyendo tendremos que:

$$z_k=z_k(1,\omega);\qquad \forall k= 1, 2, \dots, n$$

De donde será posible obtener el vector de soluciones:

$$w_k^*=w_k(\vec z)\qquad \forall k=1,2,\dots, n$$
 
Sustituyendo en la función objetivo (4):

$$\mathcal Y\equiv \{(q, \vec z)>0\;:\;q\leq w(\vec z)\cdot \vec z+\Pi(1, w(\vec z))\}$$

De acuerdo con la definición de función de beneficios, se verificará que $\Pi(1,\omega)=q-w\cdot \vec z$, por lo que:

$$\underset{(6)}{\\[1em]w(\vec z)\cdot \vec z+\Pi(1, w(\vec z))=q(w(\vec z))}$$

Y, comparando (6) y (1) vemos que:

$$f(\vec z)=q(w(\vec z))$$

Por lo que habremos recuperado la función de producción a partir de la solución del problema de minimización del coste. De manera sintetizada:

\eqblock{
\begin{aligned}
    &\hspace{0.5em}\Pi(p,w)&&\Rightarrow \hspace{2em}-\frac{\Pi(p,\vec w)}{\partial w_k}=z_k^*=z_k(p,\vec w)=z_k(1, \omega) &&\Longrightarrow \hspace{3em} w_k^*=w_k(\vec z)\\[1em]
    &\hspace{1.7em}\Downarrow&& &&\hspace{7em}\Downarrow \\[1em]
    &\frac{\partial \Pi(p, \vec w)}{\partial p}&&\Rightarrow\hspace{5.5em} q^*=q(p,\vec w)=q(1,\omega) &&\Longrightarrow\hspace{2em}\boxed{f(\vec z)=q(1, \omega(\vec z))}
\end{aligned}
}{\textcolor{red}{Poner título a esta ecuación}}
 
\subsection{Teorema básico de la dualidad}
Los axiomas anteriores nos permiten representar la tecnología mediante una función de producción, que se define como una relación puramente técnica entre insumos productivos y volúmenes de producción incluyendo todos los métodos de producción.\footnote{%
    En realidad, podemos distinguir dos categorías de axiomas:
\begin{la}
    \item Necesarios (garantizan que la tecnología pueda representarse mediante una función de producción): (i) no vacío; (ii) cerrado; (iii) eliminación gratuita; y, (iv) imposibilidad deproducción gratuita; y,
    \item Deseables (dan pie a una tecnología de buen comportamiento y, por ende, a una función neoclásica de buen comportamiento, por lo que son condiciones indispensables para la competencia perfecta): (i) posibilidad de inacción, se puede elegir si producir o no; y, (ii) convexidad, La convexidad implica la posibilidad de divisibilidad de escala (i.e. se puede reducir la escala de la producción). Además implica rendimientos a escala no crecientes. Los procesos extremos son menos productivos que los promediados. En el caso de un solo output este axioma implica que la función de producción es cóncava. Para maximizar el beneficio esta condición es esencial.
\end{la}
}

\subsection{Propiedades de la función de costes y de la demanda condicionada de factores}




Habiendo estudiado la función de producción y habiendo estudiado algunos ejemplos de funciones de producción vamos a pasar al segundo bloque de la exposición en el cual introduciremos herramientas analiticas para .la función de producción, comentando las leyes de la función de producción en el corto, largo y muy largo plazo.


\clearpage
\section{Dimensión temporal de la función de costes}
\puntosclave{
    La teoría económica distingue entre los costes a corto plazo y los costes a largo plazo. Los costes a corto plazo son aquellos en los que algún factor de producción (generalmente los activos de capital y la dirección de la empresa) son fijos.

    
	Los costes a largo plazo son aquellos a través de un periodo de tiempo lo suficientemente largo como para permitir el cambio de todos los factores de producción. En el largo plazo todos los factores son variables.\footnote{%
	    Se puede entender también del siguiente modo. Los costes de corto plazo son aquellos con los que la empresa opera en cualquier periodo. Los costes a largo plazo son costes planeados o costes ex ante en tanto que representan las posibilidades óptimas de expansión de la produccióny de esta forma ayudan al emprendedor a planear sus actividades futuras. De esta forma, previo a cualquier inversión, el empresario se encuentra en una situación de largo plazo en el sentido de que éste puede elegir una alternativa de entre un amplio conjunto de inversiones posibles definidas por el estado de la tecnología. Una vez que la decisión de inversión se ha tomado y los fondos son atados en activos de capital fijo, el empresario opera bajo condiciones de corto plazo: ahora está en una curva de coste de corto plazo.
	}

	De la representación de las curvas de costes en el corto plazo se puede observar que al mínimo del coste medio total se le denomina óptimo de explotación, ya que, si el precio es igual al mínimo, se cubren con los igresos todos los costes
    
	La curva de coste medio a largo plazo representa el coste medio mínimo para obtener cada nivel de producción, cuando se dispone del tiempo necesario para ajustar óptimamente los factores de producción y no tiene por qué coincidir con el mínimo de los costes medios a corto plazo. Al tamaño de planta para el que coincide el mínimo del coste medio a corto plazo con el mínimo del coste medio a largo plazo se le denomina dimensión óptima de la empresa o tamaño óptimo de planta
    
	Según el Teorema de LE CHÂTELIER-SAMUELSON las funciones de oferta de la empresa competitiva a corto plazo son más inelásticas que la de largo plazo y tanto más cuanto mayor es el número de factores considerados fijos.\footnote{Demostración en SEGURA (1988), pag. 196
    }
    
	Las economías de escala ocurren cuando los costes crecen menos que proporcionalmente que la producción, X, o de forma equivalente, cuando los costes medios a largo plazo son decrecientes en el nivel de producción. Los rendimientos a escala crecientes implican economías de escala, pero no al revés.
}
	
En el apartado anterior hemos visto cómo del problema de minimización de costes obtenemos dos resultados: un SCEDIC y la Función Valor de Costes.
Estos resultados se derivan de las funciones de producción que describen los métodos de producción técnicamente eficientes. Pero además, en este problema se incorporan también los precios, por lo que la función de costes garantiza además la eficiencia económica.

Supongamos dos factores productivos, capital ($K$) y trabajo ($L$). En este caso, el problema de minimización del coste de la empresa se puede especificar de la siguiente manera:

\eqblock{
\begin{aligned}
    &\boxed{
    \begin{aligned}
        &\min_{\{K,L\}}\quad C(K,L)=w_L\cdot L+w_K\cdot K\\
        &\text{s.a.}\quad
        \begin{cases}
            \bar Q\leq f(K,L)\\[0.5em]
            \{K,L\}\in \mathbb R^+
        \end{cases}
    \end{aligned}
    }\\[2em]
    &\text{SCEDIC:}\quad
    \underset{\substack{\text{\scriptsize Sustituimos en la Función Objetivo}\\\text{\scriptsize y obtenemos Función de Costes}}}{
    \begin{cases}
        L^*(\vec w, \bar Q)\\
        K^*(\vec w, \bar Q)
    \end{cases}\;\Longrightarrow}\;\boxed{C(\vec w, \bar Q)=\vec w\cdot \vec h(\vec w, \bar Q)=w_L\cdot L^*(\vec w, \bar Q)+K^*(\vec w, \bar Q)}
\end{aligned}
}{Análisis temporal de la producción: Problema de minimización. Dimensión temporal de la producción: Punto de partida}
 
A partir de la resolución del problema, obtenemos las funciones de demanda de inputs condicionada, que luego utilizamos para sustituir en la función objetivo con el fin de obtener la Función de Costes para el caso de dos factores productivos. Ahora vamos a ver tanto gráficamente como analíticamente como varían estas funciones cuando varían los parámetros de los que dependen, es decir, cuando varían los precios de los factores productivos y cuando varía la cantidad producida. Será, por tanto, análogo a lo visto en la Teoría de la demanda neoclásica del consumidor [ver Tema 3.A.9]. 

\textbf{NOTA.-} Ante variaciones en precios de los factores productivos (conceptos de elasticidad del coste respecto al precio de un factor y elasticidad del coste marginal respecto al precio de un factor, que permite clasificar entre factores normales e inferiores, págs. 221-244 del TUGORES y págs. 123 y ss del GRAVELLE) y sobre todo ante variaciones en la cantidad de output (conceptos de coste medio, marginal, elasticidad del coste [a secas, no confundir con la elasticidad del coste respecto al precio de un factor] para establecer situaciones de (des)economías de escala, y la relación entre elasticidad coste y elasticidad de escala, esto es, relación entre rendimientos de escala y economías de escala). 


\subsection{Análisis de estática comparativa}
Levantamos el supuesto de estricta cuasi-cóncavidad que impusimos al establecer las Relaciones de Dualidad y supondremos que todos los factores son inicialmente variables: largo plazo (L/p).\footnote{%
    En realidad, exponer la parte de estática comparativa o la parte de Largo Plazo es, en puridad, exactamente la misma. Esto se debe a que la estática comparativa presupone la elasticidad de todos los factores (como bien supusimos en el apartado anterior donde desarrollamos toda la Relación de Dualidad). Entonces, podemos ser inteligentes y evitar exponer las dos siempre y cuando lo liguemos de manera apropiada y entendamos bien lo que estamos haciendo.
}

\subsubsection{Variaciones en el nivel de producción (Q): Perspectiva de demanda de factores}
Veremos a continuación dos formas de analizar las variaciones sobre la demanda de factores productivos atendiendo a una variación en el nivel de producción. Estas dos perspectivas, complementarias, son fundamentalmente: el análisis gráfico y el análisis en base a elasticidades.
Trabajaremos sobre las funciones de demanda de factores que forman el SCEDIC: $\vec h^*(\vec w,\bar C)$

\paragraph{Análisis gráfico: Curva producción-demanda de factores productivos (Senda de Expansión)}
La trayectoria o senda de expansión óptima es el lugar geométrico de los puntos o combinaciones óptimas de factores que se obtienen al ir aumentando los costes y, por consiguiente la producción, manteniendo los precios de los factores constantes. 

\imagenfit{Senda_Expansion.png}{Senda de expansión. Análisis largo plazo}{Apuntes Víctor Gutiérrez. Tema 3.A.12}

A lo largo de la trayectoria de expansión óptima, la RMST es igual a la razón de los precios y, como asumimos que los precios no varían, la RMST permanece constante. Como consecuencia, la trayectoria de expansión óptima es una isoclina [\authorfont{Ver Tema 3.A.11}] a lo largo de la cual aumenta la producción, permaneciendo constante el precio de los factores de producción e indica cómo cambian las proporciones de los factores al modificarse el nivel de producción, manteniéndose constante el coste de éstos. Si la función de producción es homotética, en el caso de dos factores, el cociente de productividades marginales solo depende de la cantidad relativa de los factores, y, por tanto, dada una combinación inicial de factores, mientras se mantenga constante la relación entre el precio de los factores cualquier punto que pase por el origen y por dicha combinación pertenece a la senda de expansión.\footnote{%
    La senda de expansión se puede volver hacia atrás. En este caso hablaríamos de inferioridad del factor trabajo: elegimos menos trabajo cuando la producción aumenta por encima de un nivel de producción. De la misma manera, si la curva tuviera pendiente negativa en el sentido opuesto en algún tramo, el factor inferior sería el capital.
} En este caso, la senda de expansión es un radio que pasa por el origen de coordenadas.

\paragraph{Análisis gráfico: Curva de ENGEL del factor productivo}
De la senda de expansión óptima se puede obtener la curva de ENGEL, que relaciona la cantidad demandada de un factor de producción con el nivel de producción, es decir, nos trasladamos desde un plano cartesiano de factores a un plano cartesiano de factor-producción. A partir del análisis de la curva de ENGEL diremos que:
\begin{la}
    \item Un factor productivo es normal en coste y producción ($w,Q$) si la derivada parcial de la Función Valor de Demanda del Factor (función del factor $i$ del SCEDIC, $h_{i}$) en función de coste y producción respecto al output es mayor que cero, esto es si la demanda es creciente en la cantidad de capital (riqueza) y, por lo tanto, la pendiente de su curva de Engel es positiva;\footnote{%
        Si la función de producción es homotética, la senda de expansión es una recta de pendiente positiva y por lo tanto ambos factores productivos son normales.
    } o,
    \item Si por el contrario el efecto riqueza es negativo y su curva de ENGEL tiene pendiente negativa, se dice que el bien es inferior en ($w,Q$). Podría suceder que la pendiente dependa del tramo en el que nos encontremos (del nivel inicial de producción / riqueza)
\end{la}

\imagenfit{Curva_ENGEL_FFPPs.png}{Curva de ENGEL sobre factores productivas. Tramos de la Curva}{Apuntes Víctor Gutiérrez. Tema 3.A.12}
 
\paragraph{Análisis de elasticidades: elasticidad-producto-factor }
Además, analíticamente podemos obtener una medida de la sensibilidad de la demanda de factores productivos ante variaciones en la cantidad producida: la elasticidad de la demanda del factor productivo-cantidad producida. La elasticidad de la demanda del factor productivo respecto a la cantidad producida mide la variación porcentual en la cantidad demandada de un factor productivo ante una variación del 1\% de la cantidad producida:

\eqblock{
\begin{aligned}
    &\text{Elasticidad $\bar Q$ - factor}:&& \varepsilon_{h_i,\bar Q}=\frac{\partial h_i(\vec w, \bar Q)}{\partial \bar Q}\cdot\frac{\bar Q}{h_i(\vec w,\bar Q)}
\end{aligned}
}{Elasticidad producto factor. Dimensión temporal (largo plazo): Variaciones en el nivel de producción}

De este modo, podemos clasificar los factores productivos de la siguiente manera:
\begin{la}
    \item Factores productivos normales. Son aquellos para los que la elasticidad-producto-demanda es mayor a cero. Lo que implica que a mayor cantidad a producir, mayor cantidad demandada del factor productivo;
    \item Factores productivos frontera. Son aquellos para los que su elasticidad-producto-demanda. Lo que implica que si varía la cantidad a producir, la cantidad demandada del factor productivo no variará; y,
    \item Factores productivos inferiores. Son aquellos para los que su elasticidad-producto-demanda es inferior a cero. Lo que implica que a mayor cantidad a producir, menor cantidad demandada del factor productivo.
\end{la}

Como es normal, la elasticidad podría depender del nivel de producción de partida, y por lo tanto un factor productivo puede ser normal para ciertos niveles de producción e inferior para otros.

\subsubsection{Variaciones en el nivel de producción ($Q$): Perspectiva de costes}
\textbf{NOTA.-} Entender que lo que estamos haciendo es en el mismo gráfico que antes, requerir una curva isocuanta más alejada del origen de coordenadas, por lo que el coste total aumenta a medida que aumenta $\bar Q$, pero la forma que adopte la función de Coste Total (cóncava o convexa) dependerá de la forma funcional de la función de producción (rendimientos crecientes o decrecientes a escala). Hacer aquí todo el análisis de largo plazo, hablar de coste marginal y coste medio, de economías de escala (internas y externas) y hacer el análisis de elasticidades. Luego extrapolar el análisis a corto plazo en el apartado de análisis de plazos.
Habiendo visto cómo varían las demandas de factores productivos ante cambios en la cantidad producida, pasamos a realizar el mismo análisis para la función de costes.
 
\paragraph{Análisis gráfico: las Curvas de Costes, $CT(Q)$, $CMe(Q)$, $CMg(Q)$}
Si suponemos un vector de costes constante, podemos representar la Función Valor de Costes solo en función del output C(Q) sobre el plano cartesiano.\footnote{%
    Supondremos que se cumple (para esta función de producción) la Ley de proporciones variables, según la cual -menteniendo un nivel fijo de un factor productivo- si incrementamos la cantidad del otro factor productiva, llegará un momento en el que su PMg disminuirá, de tal forma que llegue incluso a ser negativa a partir de un determinado punto. Esto da lugar al tipo de Gráfico que vemos (S invertida), y nos es conveniente porque da lugar a todo tipo de rendimientos y además da lugar a diferentes tramos de economías de escala (alguno de economías de escala y otros de deseconomías de escala). Sin embargo, si la función deproducción es neoclásica de buen comportamiento y convexa en todos sus tramos (p.e. tipo Cobb-Douglas), no se dará este caso. Esta Ley también se conoce como la Ley de Rendimientos Marginales decrecientes.
}

\imagenfit{CurvaCostes_VarProduccion.png}{Curva de costes para un mismo nivel de producción}{Apuntes Víctor Gutiérrez. Tema 3.A.12}

Los Costes Medios ($CMe$) se definen como los $CT(Q)$ divididos entre la cantidad producida (o como la pendiente de la secante que parte del origen de coordenadas).\footnote{%
    Para representar el coste medio hay que fijarse en la pendiente de la secante que parte desde el origen de coordenadas. Los CMe decrecerán hasta que la recta secante a la curva de Costes Totales sea tangente a la curva Costes Totales y a partir de ese punto comienza a crecer.
} También podemos definir los costes marginales como el coste de producir una cantidad marginalmente mayor y se calcula como la pendiente de los costes totales en cada punto (o como la pendiente de la tangente en cada punto).\footnote{%
    Para representar el coste marginal hay que fijarse en la curvatura de la función de costes totales. Los CMg decrecerán cuando la función Costes Totales sea cóncava (la pendiente decrece a medida que aumenta la producción) y crecerá cuando la función Costes Totales sea convexa. El mínimo del coste marginal se da en el punto de inflexión de la curva de costes totales.
}

\eqblock{
\begin{aligned}
    CMe=\frac{C(\vec w,\bar Q)}{\bar Q}\qquad \text{y}\qquad CMg=\frac{\partial C(\vec w, \bar Q)}{\partial \bar Q}
\end{aligned}
}{Definición Costes Medios y Costes Marginales. Dimensión temporal de la producción: Variaciones en el nivel de proónducci}

La curva de $CMe(Q)$ a la que da lugar esta función de Costes Totales tendrá, a largo plazo, forma de U.\footnote{%
    WICKSELL fue el primer economista que señaló que podrían existir rendimientos mixtos a escala, es decir, para todo nivel de producción es posible que la empresa pudiera pasar por los tres rendimientos: al principio de la expansión serían crecientes, pero pasarían a ser decrecientes y en el punto de inflexión serían constante.
} Como el coste del factor se mantiene idéntico para todos los niveles de producción, cuando la PMg comience a decrecer, afectará negativamente a los costes (los aumentará) y a partir de ese punto la función de Costes Totales será convexa, es decir, crece aceleradamente.\footnote{%
    En este caso estamos suponiendo que dichos rendimientos marginales decrecientes aparecen en un cierto punto, es decir, la PMg del factor elástico / variable aumenta hasta un cierto punto, a partir del cual al aumentar la cantidad del factor variable su PMg comienza a reducirse. Ello es lo que origina que el CMe aumente a partir de un determinado punto.
}\textsuperscript{,}\footnote{%
    Esta es la clave de que los rendimientos marginales crecientes (PMg creciente) implique economías de escala pero que las economías de escala no necesariamente impliquen rendimientos marginales crecientes (PMg puede ser creciente o decrecientes).
}

\imagenfit{CurvasCostes.png}{Curvas de costes totales, medios y marginales. Proyección y propiedades}{Apuntes Víctor Gutiérrez. Tema 3.A.12}

La curva de CMe será decreciente hasta el punto donde se igualan con los CMg. A partir de este punto, será creciente, teniendo por tanto forma de U:
\begin{la}
    \item En el tramos de CMe decrecientes, al aumentar el nivel de producción los costes aumentan menos que proporcionalmente, por lo tanto diremos que estamos ante economías de escala; y, 
    \item Por otro lado, en el tramos en que los CMe son crecientes, los costes aumentarán proporcionalmente más que la producción, por lo que diremos que estamos ante deseconomías de escala.
\end{la}

\paragraph{Análisis de elasticidades: elasticidad-coste-producto}
Podemos medir el concepto de las economías de escala mediante la elasticidad del coste respecto a la cantidad producida. La elasticidad del coste respecto a la cantidad producida mide la variación porcentual en el coste ante una variación porcentual de la cantidad producida:

\eqblock{
\begin{aligned}
    &\varepsilon_{C,\bar Q}=\overbrace{\frac{\partial C(\vec w, \bar Q)}{\partial \bar Q}}^{\equiv CMg(\bar Q)}\cdot \underbrace{\frac{\bar Q}{C(\vec w, \bar Q)}}_{\equiv\frac{1}{CMe(\bar Q)}}=\frac{CMg(\bar Q)}{CMe(\bar Q)}\\[2em]
    &\text{y, si}\quad 
    \left\{\begin{aligned}
        &\varepsilon_{C,\bar Q}<1\Longrightarrow &&\underset{\text{\scriptsize $\Rightarrow CMe$ decrecientes}}{CMg(\bar Q)<CMe(\bar Q)}\Longrightarrow&&\text{Economías de escala}\\[1em]
        &\varepsilon_{C,\bar Q}=1\Longrightarrow &&\underset{\text{\scriptsize $\Rightarrow CMe$ constantes}}{CMg(\bar Q)=CMe(\bar Q)}\Longrightarrow&&\text{Ausencia economías de escala}\\[1em]
        &\varepsilon_{C,\bar Q}>1\Longrightarrow &&\underset{\text{\scriptsize $\Rightarrow CMe$ crecientes}}{CMg(\bar Q)>CMe(\bar Q)}\Longrightarrow&&\text{Deseconomías de escala}
    \end{aligned}\right.
\end{aligned}
}{Elasticidad coste producto. Dimensión temporal (largo plazo): Variaciones en el nivel de producción}
 
De hecho, podemos ver, con ayuda del Gráfico de Costes Totales como una misma función puede presentar economías mixtas (la Función de Producción presenta zonas de concavidad, convexidad y escalabilidad constante, rendimientos mixtos). En el sentido de que puede iniciar presentando economías de escala, tener un periodo con ausencia de economías de escala y luego presentar deseconomías de escala.

\imagenfit{RendMixtos_Escala.png}{Rendimientos mixtos a escala. Fases de producción}{Apuntes Víctor Gutiérrez. Tema 3.A.12}

\paragraph{Relaciones: elasticidad-escala y elasticidad-coste-producto}
Aunque resulta tentador, el concepto de economías de escala [Tema 3.A.12] no puede ser confundido con el concepto de rendimientos de escala [\authorfont{Tema 3.A.11}].\footnote{%
    La elasticidad-escala mide la variación porcentual en el nivel de producción ante una variación del porcentual equiproporcional de todos los factores productivos:
$$\varepsilon_{Q,\lambda}=\frac{\mathrm dQ}{\mathrm d\lambda}\cdot\frac{\lambda}{Q}=\underset{\varepsilon_{Q,L}}{\frac{{PMg}_L}{{PMe}_L}}+\underset{\varepsilon_{Q,K}}{\frac{{PMg}_K}{{PMe}_K}}$$
} Las economías de escala son un concepto más amplio que los rendimientos a escala, dado que no exigen que la variación que experimenta la producción tenga su causa en una variación equiproporcional de los factores productiva (salto de isocuantas).

De esta forma, aunque resulta tentador pensar que las economías de escala surgirán únicamente cuando la función de producción exhiba rendimientos crecientes a escala y que habrá deseconomías de escala cuando la función de producción exhiba rendimientos decrecientes a escala esto no siempre será así:
\begin{lnum}
    \item En el corto plazo, podría darse el caso que una función de producción, la función de costes exhibiese economías de escala incluso cuando la función de producción exhibiese rendimientos a escala decrecientes pues la empresa puede obtener economías de escala con el aumento de un solo factor productivo (si así se experimenta);\footnote{%
        Por ejemplo, imaginémos que los precios de los factores no son constantes (por descuentos, aprendizaje, indivisibilidades, mejor utilización de equipos fijos, etc.), la empresa puede lograr economías de escala en costes aunque su tecnología tenga rendimientos decrecientes a escala.
    }
    \item En el largo plazo, rendimientos decrecientes sí implican deseconomías de escala salvo que se rompa alguna de las hipótesis estándar del largo plazo. Es decir, sí podría suceder, pero solo si:
    \begin{la}
        \item Los precios de los factores no son constantes (por descuentos, poder de mercado o aprendizaje); o,
        \item Existen eficiencias externas o tecnológicas que reducen los costes al crecer la escala (por ejemplo, economías externas o de red). En ese caso, la función de costes puede mostrar economías de escala aparentes, aunque la tecnología interna tenga RDE.
    \end{la}
    Pero eso ya no proviene de la tecnología de producción, sino de factores externos o de mercado.
\end{lnum}

En cambio, si la función de producción es homotética, la relación entre economías de escala y rendimientos a escala si será biunívoca. Concretamente, teniendo en cuenta que los rendimientos a escala de una función de producción se pueden medir mediante la elasticidad-escala, la relación entre ambos conceptos en el caso de funciones homotéticas vendrá dada por:

\eqblock{
\begin{aligned}
    \underset{\text{\scriptsize Función producción homotética}}{\text{Caso particular:}}&&\varepsilon_{C,\bar Q}=\frac{1}{\varepsilon_{Q,\lambda}}
\end{aligned}
}{Relación: Elasticidad-escala y Rendimientos a escala. Dimensión temporal: Variaciones en el nivel de producción}
 
\textbf{NOTA.-} Por ejemplo, cuando el factor productivo variable (en nuestro caso L) exhibe rendimientos marginales crecientes los costes marginales son decrecientes. Del mismo modo, cuando los rendimientos marginales son decrecientes, los costes marginales serán crecientes, formándose una curva de costes marginales en forma de U.

\subsubsection{Variaciones en el coste de los inputs ($\vec w$): Perspectiva de demanda de factores}
Veremos ahora como varía la demanda de factores productivos cuando experimentan una variación en su precio (o coste).

\paragraph{Análisis gráfico: Curva precio-demanda de factores productivos}
\textbf{NOTA.-} Aquí se está representando el problema Primal de SEGURA, lo óptimo sería cambiar las gráficas de manera que la restricción fuera el nivel de producción y cambiara el precio de los bienes. Esto daría pie a que no puede existir factores productivos Gifenn

Si modificamos el precio de un factor de producción (p. ej. trabajo), las curvas isocoste rotan sobre el punto en el que estas cortan el eje del bien cuyo precio permanece constante. La curva precio-demanda de factores productivos une los puntos de tangencia entre las curvas isocoste y las curvas isocuantas a medida que aumenta el precio de uno de los dos factores.

\imagenfit{Curva_Precio_Demanda_Factores.png}{Curva de precio-demanda de factores productivos}{Apuntes Víctor Gutiérrez. Tema 3.A.12}
 
\paragraph{Análisis gráfico: Curva inversa de demanda del factor productivo (en función de su coste)}
Representando cantidad del factor demandada para cada nivel de precio, obtenemos la curva inversa de demanda del factor trabajo:

\imagenfit{CurvaInversa_DemandaFP.png}{Curva inversa de demanda del factor productivo}{Apuntes Víctor Gutiérrez. Tema 3.A.12}
 
De hecho, es la construcción inversa de la curva anterior.

\paragraph{Análisis de elasticidades: elasticidad-precio-propio}
Analíticamente podemos obtener una medida de la sensibilidad de la demanda ante variaciones en los precios: la elasticidad de la demanda del factor productivo-precio propio y la elasticidad de la demanda del factor productivo-precio cruzado. La elasticidad de la demanda del factor productivo respecto al precio propio mide la variación porcentual en la cantidad demandada de un factor productivo ante una variación del 1\% de su precio:

\eqblock{
\begin{aligned}
    \text{Elasticidad precio-demanda}:&&\varepsilon_{h_i,w_i}=\frac{\partial h_i(\vec w, \bar Q)}{\partial w_i}\cdot\frac{w_i}{h_i(\vec w,\bar W)}
\end{aligned}
}{Elasticidad de la demanda del factor productivo-precio propio. Dimensión temporal de la producción: Variaciones en precios}
 
Recordemos que, en el caso de la Teoría neoclásica de la demanda del consumidor, si la elasticidad precio-propio era superior a cero decíamos que este podía ser un bien Giffen o Veblen (bien inferior con supuesto adicional ER > ES). Aquí esto no puede suceder ya que la Función de Costes tiene como propiedad que su matriz de sustitución es semidefinida-negativa, lo que implica que (haciendo uso del Lema de Shepard) la segunda derivada de la función de costes respecto al coste de un factor es igual a la derivada de la función de demanda respecto al precio de ese factor y, ambas, menores o iguales a cero. Por lo tanto se cumple la Ley de la Demanda Compensada.\footnote{%
    Las curvas de demanda marshalliana con pendiente descendente muestran el efecto de los cambios en los precios sobre la cantidad demandada. A medida que sube el precio de un bien, la cantidad del bien demandado disminuirá, manteniendo la riqueza y otros precios constantes.
}

\paragraph{Ecuación de SLUTSKY para Teoría de la Producción \textcolor{red}{OJO: Encajar en el apartado anterior}}
Si existe un trasvase correcto y completo de información en la Relación Dual sabemos que los costes en el PMaxP coinciden con el coste mínimo en el PMinC para el nivel de producción que ha sido maximizado, entonces:

\eqblock{
\begin{aligned}
&\text{Ecuación de SLUTSKY: Teoría de la Producción}\\[1em]
    &\over\begin{aligned}
        \\[0.5em]
        &\text{Principio de dualidad:}\hspace{5em}\boxed{h_i(\vec w, \bar Q)=z_i(\vec w, C(\vec w, \bar Q))}\\[3em]
        &\left\{\begin{aligned}
            &\underset{\substack{\text{\scriptsize Derivamos respecto}\\\text{\scriptsize a $w_j$ cualquiera}}}{\Longrightarrow}
            &&\begin{aligned}
                &\frac{\partial h_i(\vec w, \bar Q)}{\partial w_j}=\frac{\partial z_i(\vec w, C(\vec w, \bar Q))}{\partial w_j}\\[0.5em]
                &\frac{\partial h_i(\vec w,\bar Q)}{\partial w_j}=\frac{\partial z_i(\vec w,\bar C)}{\partial x_j}+\frac{\partial z_i(\vec w,\bar C)}{\partial w_j}\cdot\left(\frac{\partial C(\vec w,\bar Q)}{\partial w_j}\right|_{\substack{=z_j^*\\\text{\scriptsize Por el Lema}\\[0.2em]\text{\scriptsize de SHEPARD}}}
            \end{aligned}
        \end{aligned}\right\}\\[3em]
        &\text{SLUTSKY:}\;\Longrightarrow\hspace{5em}\boxed{\underbrace{\frac{\partial z_i(\vec w,\bar C)}{\partial w_j}}_{\text{\scriptsize Efecto Total}}=\overbrace{\frac{\partial h_i(\vec w,\bar Q)}{\partial w_j}}^{\text{\scriptsize Efecto Sustitución}}-\underbrace{z_j\cdot \frac{\partial z_i(\vec w,\bar C)}{\partial w_j}}_{\text{\scriptsize Efecto Renta}}}
    \end{aligned}
\end{aligned}
}{Ecuación de SLUTSKY en Teoría de la Producción: Principio de dualidad. Dimensión temporal: variaciones en los precios}
 
Por tanto, dadas las circunstncias correctas, derivando la relación/principio de dualidad, obtenemos la ecuación de SLUTSKY análoga al caso del problema del consumidor [ver Tema 3.A.9]. Para analizar el efecto sobre cambios en el precio del mismo factor sobre su propia demanda, analizamos para el caso de $i = j$. Esto nos permite descomponer el efecto de cambios en los precios de un factor productivo en dos efectos contrapuestos: el efecto sustitución propio y el efecto renta propio.

\paragraph{Análisis de elasticidades: elasticidad-precio-cruzado}
La elasticidad de la demanda del factor productivo respecto al precio cruzado mide la variación porcentual en la cantidad demandada de un factor productivo ante una variación del 1\% del precio de otro:

\eqblock{
\begin{aligned}
    \varepsilon_{h_i,w_j}=\frac{\partial h_i(\vec w,\bar Q)}{\partial w_j}\cdot\frac{w_j}{h_i(\vec w,\bar Q)}
\end{aligned}
}{Elasticidad de la demanda del factor productivo-precio cruzado. Dimensión temporal de la producción: Variaciones en precios}
 
Con la ecuación de SLUTSKY podremos descomponer los efectos de cambios en el coste de un factor productivo sobre la cantidad demandada de otro y analizar las relaciones de sustituibilidad y complementariedad brutas entre los factores productivos.

\subsubsection{Variaciones en el coste de los inputs ($\vec w$): Perspectiva de costes}
Habiendo visto como varían las demandas de factores productivos ante cambios en sus precios, pasamos a realizar el mismo análisis para la función de costes:

$$C(\vec w, \bar Q)=w_L\cdot L^*(\vec w,\bar Q)+w_K\cdot K^*(\vec w, \bar Q)$$
 
El efecto que tendrá un aumento proporcional en el coste de los inputs sobre el coste total de la empresa es directo. Esto es así dado que, como ya hemos dicho, la función de costes es homogénea de grado uno en coste de los factores:

$$C^*(\theta\cdot\vec w,\bar Q)=\theta^1\cdot C^*(\vec w, \bar Q)=\theta C^*(\vec w,\bar Q)$$
 
Los efectos que tiene el cambio del precio de un solo factor sobre el coste total también son bastante directos. Esto lo podemos ver tanto analíticamente como gráficamente. Desde un punto de vista gráfico, variaciones en el precio de los factores productivos se traducirán en desplazamientos de la curva de costes totales en el plano ($C,\bar Q$), aumentando del mismo modo los costes medios para todo nivel de producción $\bar Q$.

\paragraph{Análisis de elasticidades: elasticidad-precio-propio}
Además, aplicando el lema de Shephard, la elasticidad del coste respecto al precio propio se puede escribir como:

$$\varepsilon_{C,w_i}=\frac{\partial C(\vec w,\bar Q)}{\partial w_i}\cdot \frac{w_i}{C(\vec w,\bar Q)}=\frac{h_i(\vec w,\bar Q)\cdot w_i}{C(\vec w,\bar Q)}$$
 
La sensibilidad del coste a un cambio en el precio de un factor individual es igual a la proporción que representa el gasto de ese factor sobre el coste total.\footnote{%
    Dado que el coste medio se mantiene constante, para determinar el efecto de la variación del coste sobre el coste medio sería necesario calcular la elasticidad del coste medio respecto a al coste del factor ($i$) que también es igual a la participación en el gasto del factor $i$:
$$\varepsilon_{{CMe}_i,w_i}=\frac{\partial C(\vec w,\bar Q)/\bar Q}{\partial w_i}\cdot\frac{w_i}{C(\vec w,\bar Q)/\bar Q}=\frac{h_i(\vec w,\bar Q)}{\bar Q}\cdot\frac{w_i}{C(\vec w_i,\bar Q)/\bar Q}=\frac{h_i(\vec w,\bar Q)\cdot w_i}{C(\vec w,\bar Q)}$$
El efecto de un aumento dado en el coste del factor ($w_{i}$) sobre las curvas de costes totales y costes medios será desplazar estas curvas verticalmente hacia arriba en una cantidad que dependerá de la proporción que suponga el gasto en el factor productivo $i$ sobre el gasto total. Esto no significa que las curvas se desplacen en la misma proporción para todos los niveles de producción, dado que es posible que la proporción de lo que supone el gasto en ese factor productivo no sea constante con el nivel de producción.
}

\subsection{Corto plazo y Largo plazo}

\subsubsection{Corto plazo: C/p}
Empezaremos trabajando sobre el corto plazo pues debemos resolver un problema de optimización adicional para trabajar sobre este. Consiste en restringir lo expuesto inicialmente para las relaciones de dualidad al supuesto de que solo un factor productivo es variable (o varios, lo importante es que haya algunos fijos y otros variables). Esto es la esencia del corto plazo.


\paragraph{Desarrollo}
Entonces, para dos factores productivos, suponiendo que el factor capital ($K$) es fijo y el factor trabajo ($L$) es variable, tendríamos:

$$\begin{aligned}
    \min_{L}\;C(L)=\overbrace{w_L\cdot L}^{\text{\scriptsize Costes variables}}+\underbrace{w_K\cdot K}_{\text{\scriptsize Costes fijos}}\\[1em]
    \text{s.a.}\quad
    \left\{\begin{aligned}
    &\bar Q\leq f(\bar K,L)\\[0.5em]
    &0\leq L
    \end{aligned}\right.
\end{aligned}$$

Vemos que el problema de minimización del coste a corto plazo es idéntico al problema general de largo plazo descrito anteriormente, salvo por el hecho de incorporar una restricción adicional correspondiente a la limitación de los factores fijos (i.e. tamaño de planta). Para resolver este problema tendremos que elegir la cantidad de trabajo óptima para una determinada producción, es decir, la demanda condicionada del factor trabajo.
En el óptimo, la empresa utilizará el capital fijado y contratará trabajo de acuerdo a una función de demanda condicionada de trabajo. También podremos expresar la función de valor del problema (función de costes) como la suma de los costes totales a corto plazo (i.e. suma de los costes variables y los costes fijos):

$$\begin{aligned}
    L^*(\vec w,\bar K,\bar Q)\;\Longrightarrow\;\boxed{C^*(\vec w,\bar Q,\bar K)=w_L\cdot L^*(\vec w,\bar Q,\bar K)+w_K\cdot K}
\end{aligned}$$
 
\paragraph{Costes Totales: C/p vs. L/p}
Por la limitación del factor fijo, los costes de corto plazo asociados a la producción de un cierto nivel de output no pueden ser menores que los correspondientes costes de largo plazo. El gráfico de la senda de expansión permite comprobar como a corto plazo los costes van a ser mayores o iguales que a largo plazo.

\imagenfit{CT_Produccion_Senda.png}{Comparativa corto y largo plazo. Costes de producción}{Apuntes Víctor Gutiérrez. Tema 3.A.12}

La senda de expansión óptima (objetivo al que llegar a corto plazo) es el lugar geométrico de los puntos o combinaciones óptimas de factores productivos que se obtienen al ir incrementando la producción manteniendo los precios de los factores constantes (por combinaciones óptimas, nos referimos a la combinación óptima de largo plazo cuando todos los factores son variables, es decir, la tangencia entre isocoste e isocuanta).
Para producir la primera isocuanta ($\bar Q_{1}$), lo ideal es situarse en el punto $A$. Pero si el capital está fijo en $\bar K$, para producir esa misma cantidad nos tenemos que situar en una isocoste mayor, concretamente en el punto $B$. En cambio, para producir la segunda isocuanta ($\bar Q_{2}$) lo ideal es situarse en el punto $C$ cuya dotación de capital coincide con el capital fija, por lo que los costes a corto plazo son iguales que los costes a largo plazo y la asignación de capital es óptima: es decir, estamos en el tamaño de planta óptimo.

En conclusión, como la empresa a corto plazo no puede redimensionarse en todos los factores para elegir el tamaño de planta óptimo conjunto, la solución al problema de optimización a corto plazo coincidirá con la solución al problema a largo plazo, pues se situará en la dotación de factores que forme parte del conjunto de soluciones a L/p.

De este modo, la empresa a corto plazo siempre producirá de forma técnicamente eficiente y a largo plazo producirá también de forma económicamente eficiente. Sin embargo, a corto plazo no podemos asegurar que esté produciendo de forma ecnómicamente eficiente, ya que sólo será así si el nivel del factor productivo fijo coincide con el nivel de factor productivo que contrataría si pudiera elegir qué cantidad de este usar. 

\paragraph{Representación gráfica: Desagregación de costes a corto plazo}
Al analizar los costes en el corto plazo, lo primero que se debe saber es que existen costes fijos y costes variables:
\begin{la}
    \item \textbf{Costes Fijos}. Incluyen principalmente factores que requieren tiempo para cambiar (p.ej. el alquiler, la tecnología, etc.), y a menudo están asociados con los costes indirectos de producción (todo lo que no es un factor de producción directo). No dependen del nivel de producción y no se pueden ajustar de acuerdo a este; y,
    \item \textbf{Costes Variables}. Dependen del nivel de producción y pueden incluir costes asociados a materias primas, la mano de obra (sin incluir salarios fijos), etc.\footnote{%
        Sólo se incluirán aquellos salarios que guarden relación con la cantidad producida o las horas trabajadas para realizar producto (output)
    }
\end{la}
	
Analíticamente, habíamos obtenido que la función de costes a corto plazo sería:

\eqblock{
\begin{aligned}
    C^*(\vec w, \bar Q,\bar K)=w_L\cdot L^*(\vec w, \bar Q,\bar K)+w_K\cdot \bar K
\end{aligned}
}{}

Los $CV(Q)$ son curvados en forma de S: (i) en un primer momento, crecen cada vez más lentamente a medida que aumenta la producción debido a los rendimientos marginales crecientes; (ii) llegado cierto nivel, alcanzan un punto de inflexión; y, (iii) comienzan a crecer aceleradamente debido a los rendimientos marginales decrecientes.

Podemos hallar los costes marginales como la pendiente de los costes totales en cada punto (o como la pendiente de la tangente en cada punto) y los CMe como los $CT(Q)$ divididos entre la cantidad producida (o como la pendiente de la secante que parte del origen de coordenadas).\footnote{%
    Para representar el coste marginal hay que fijarse en la curvatura de la función de CT(Q). Los CMg decrecerán cuando la función CT sea cóncava (la pendiente decrece a medida que aumenta la producción) y crecerá cuando la función CT sea convexa. El mínimo del coste marginal se da en el punto de inflexión de la curva de costes totales.
}\textsuperscript{,}\footnote{%
    Para representar el coste medio hay que fijarse en la pendiente de la secante que parte desde el origen de coordenadas. Los CMe decrecerán hasta que la recta secante a la curva de CT sea tangente a la curva CT y a partir de ese punto comienza a crecer.
} La forma en que costes fijos y variables afectan a la producción está relacionada con las esconomías a escala. Los rendimientos a escala se estudian analíticamente mediante la elasticidad-producción-coste (se trata de la escala económica):

\eqblock{
\begin{aligned}
    \varepsilon_{C,\bar Q}=\underset{\equiv CMg}{\frac{\partial C(\vec w,\bar Q)}{\partial \bar Q}}\cdot \underset{\equiv\frac{1}{CMe(Q)}}{\frac{\bar Q}{C(\vec w, \bar Q)}}=\frac{CMg(Q)}{CMe(Q)}
\end{aligned}
}{Elasticidad producción-coste. Dimensión temporal: C/p vs. L/p}

Vamos a verlo.

\paragraph{Representación gráfica: Costes y Economías a escala}
Dividiremos el eje de abcisas en tres Fases, tal y como hicimos para el estudio de la $PMg$ y la $PMe$, con el fin de averiguar que rendimientos existen en cada Fase.

Sobre el plano cartesiano trazamos las lineas de $CT(Q)$, $CV(Q)$ y $CF$. Obtenemos las derivadas de la función: $CMg(Q)$; y también $CTMe(Q)$. Además, podemos también desagregar los $CTMe(Q)$ en dos subtipos: $CVMe(Q)$ y $CFMe(Q)$.
\begin{la}
    \item \textbf{FASE I}. Donde la elasticidad-producto-coste es menor a uno ($CMg<CMe$), hay economías de escala. Esto implica que el coste total varía en menor proporción que la producción (i.e. coste medio decreciente). Esta Fase iría desde el origen hasta el final del punto de inflexión experimentado por la función de $CV(Q)$, es decir, cuando los $CMg$ empiezan a crecer sustancialmente rápdo, $CMg > CV$ (y $CMg < CVMe$);
    \item \textbf{FASE II}. En esta Fase (puntual) la elasticidad-producto-coste es igual a uno ($CMg=CMe$). Esto implica que el coste total varía en la misma proporción que la producción (no hay economía a escala); y,
    \item \textbf{FASE III}. Que se inicia cuando el $CMg > CVMe$, existen deseconomías de escala. Esto implica que el coste total varía en mayor proporción que la producción (i.e. coste medio creciente).
\end{la}

\imagenfit{Analisis_FASES.png}{Análisis por Fases de CASSEL. Corto y largo plazo}{Apuntes Víctor Gutiérrez. Tema 3.A.12}

\subsubsection{Corto plazo (C/p) vs. Largo plazo (L/p): Principio de LE CHATELIER-SAMUELSON}
En el largo plazo, la curva de costes totales es la envolvente de todas las curvas de costes totales a corto plazo. Del mismo modo, la curva de costes medios es la envolvente de todas las curvas de costes medios a corto plazo.
En relación a las curvas de $CMg$ a C/p, éstas cortan a las curvas de $CMe$ a C/p por su mínimo. Además, las curvas de $CMg$ a C/p cortan a la curva de $CMg$’s a L/p en el nivel de producción en el que la curva de $CMe$ a C/p es tangente a la $CMe$ a L/p.

\imagenfit{Ppio_CHATELIER_SAMUELSON.png}{Principio de LE CHATELIER-SAMUELSON}{Apuntes Víctor Gutiérrez. Tema 3.A.12}

Nótese que para el nivel de producción en que el CT a C/p es tangente con el CT a L/p, la curva de CMe a C/p es tangente a la curva de CMe a L/p y además, la curva de CMg a C/p corta a la curva de CMg a L/p.

Otra cuestión que observamos a partir del gráfico es que el tramo creciente de las curvas de CMg es más rígido en el caso C/p que el caso de L/p. Esto es relevante porque en competencia perfecta, el tramo creciente de la curva de CMg a partir del umbral de producción es precisamente la curva de oferta. El fenómeno de la mayor rigidez se explica por el Teorema de LE CHÂTELIER-SAMUELSON, de acuerdo con el cual las funciones de oferta de la empresa precio-aceptante en el C/p son más rígidas que en el L/p (y tanto más rígidas cuanto mayor sea el número de factores considerados fijos) debido a que la capacidad de adaptación de la empresa en el corto plazo es menor al tener algunas dotaciones fijas.\footnote{%
    SAMUELSON es especialmente conocido por el planteamiento general del método de las estáticas comparativas que hizo en su libro Fundamentos del análisis económico (1947). En este trabajo, SAMUELSON comenzó a generalizar la aplicación de modelos matemáticos a la economía. En particular, los estudios de termodinámica de WILLARD GIBBS (con base en la evolución de los principios del químico francés LE CHÂTELIER) llevaron a SAMUELSON a establecer el método de estática comparativa en economía. Este método explica los cambios en la solución de equilibrio de un problema de maximización obligada cuando una de las restricciones es marginalmente reforzada o relajada. Esta extrapolación del estudio de los principios termodinámicos ha servido para desarrollar los distintos escenarios económicos que se originan cuando se altera alguna variable del sistema estudiado.
Las aplicaciones matemáticas en economía y sus nexos con la química, le valieron para que en 1970 fuera galardonado con el Premio Nobel de Economía. Obtuvo esta distinción por estos estudios, en los que introdujo una aplicación sistemática de la metodología de maximización a un amplio conjunto de problemas. La introducción matemática de las teorías keynesianas con la corriente económica de la síntesis neoclásica se le debe sin duda a SAMUELSON. Su enfoque matemático y su liderazgo intelectual llevó al estudio de la economía al camino del análisis matemático que, desde entonces ha dominado la profesión de las económicas.
}

\subsection{Extensiones a la Teoría neoclásica de los Costes}


\subsection{Críticas a la Teoría Neoclásica de los Costes}
\puntosclave{
    STIGLER (1939) custiona la forma en U de los costes medios a largo plazo de la teoría tradicional, proponiendo costes medios variables constantes, con untramo horizontal, para los niveles de producción usuales (reserva de cpacidad), empezando a crecer si la empresa trabaja sobrecargada y utiliza demasiado factor variable en relación con el factor fijo o la capacidad de la empresa

	LIEBESTEIN argumenta que, en condiciones normales, dada la existencia generalizada de mercados de competencia imperfecta, el comportamiento minimizador del coste es una excepción. Se produce, por el contrario, la llamada ineficiencia X, o despilfarro de recursos ocasionada por la falta de competencia externa a la empresa, un incentivo a producir ineficientemente. 

	Existen economías de alcance cuando el coste mínimo de producir los bienes por separado supera al de hacerlo conjnutamente
	Una tecnología de producción presenta subaditividad estricta en costes si resulta menos costoso producir un output dado en una sola empresa que repartiéndolo entre varias
    
	La empresa puede ir reduciendo costes gracias a las economías debidas al aprendizaje del trabajador y de la gestión, fenómeno que se produce conforme crece la producción acumulada en el tiempo
}

La teoría neoclásica, tal y como ha sido planteada a lo largo de esta exposición hasta este punto, ha sido duramente cuestionada tanto respecto a sus fundamentos teóricos como a sus implicaciones empíricas. Las críticas han partido esencialmente de los autores de la Teoría de la Organización Industrial.

\subsubsection{La Teoría de Costes de STIGLER: la reserva de capacidad}
GEORGE STIGLER (1939) cuestiona la forma en U de los costes medios de la teoría neoclásica, tanto a nivel empírico como teórico. Considera que como existen factores productivos que no son perfectamente divisibles y existe incertidumbre sobre cómo va a ser la demanda, la función de costes va a tener un tramo en el que los costes medios y los costes marginales son constantes. Como consecuencia, no habría un tamaño óptimo, sino que habría una horquilla de tamaños óptimos que el denomina: reserva de capacidad.

\imagenfit{STIGLER_ReservaCapacidad.png}{Reserva de capacidad: STIGLER}{Apuntes Víctor Gutiérrez. Tema 3.A.12}
 
STIGLER argumenta que esto se debe, principalmente a dos razones:
\begin{lnum}
    \item \textbf{Razones económicas}. Por ejemplo, posibles cambios en la demanda del bien debido a razones cíclicas, cambios favorables en los gustos de los consumidores, o como un método preventivo para evitar una parada de la producción en caso de que haya una avería o una reparación.
    \item \textbf{Razones técnicas}. Por ejemplo, el hecho de algunas instalaciones básicas son indivisibles, o que gerentes y trabajadores son contratados en exceso para permitir futuras ampliaciones.
\end{lnum}
	
Hay que señalar que esta teoría permite que la máxima eficiencia exista durante cierto tiempo, lo que difiere completamente de la visión neoclásica, que ve la máxima eficacia como un punto de inflexión, en lugar de una línea recta.

\imagenfit{CMe_LP_STIGLER_ReservaCapacidad.png}{Curva de costes medios a largo plazo. Teoría de Reserva de capacidad. STIGLER}{Apuntes Víctor Gutiérrez. Tema 3.A.12}

Por otro lado, a largo plazo tendríamos una curva de $CMe$ tal que, los costes de producción disminuyen debido a economías de escala hasta que se alcanza el tamaño óptimo, y luego permanecen constante. Una mayor reducción del coste medio de producción podría darse gracias a una mejora de la tecnología o de las habilidades de los trabajadores. En cuanto a los costes de gestión, STIGLER considera que empiezan a aumentar después de que se alcance un cierto tamaño. Sin embargo, las deseconomías de escala de gestión serían compensadas por las economías de escala técnicas, con un efecto neto de suma cero y por lo tanto se mantendrían constantes.

\subsubsection{LEIBENSTEIN: Ineficiencia X}
La teoría anterior cuestiona la forma de U de las curvas de coste medio. Sin embargo, no pone en duda cuestiones fundamentales como la minimización del coste por parte de la empresa, que es precisamente lo que hace HARVEY LEIBENSTEIN (1966).
Una ineficiencia X surge como resultado de que ciertos inputs no den lugar a la mayor cantidad de output correspondiente como consecuencia de cierto factor X. Esto se traduce en un fallo de minimización de costes y maximización de producción e implica una pérdida de eficiencia. Este autor pone en duda el óptimo del productor, partiendo de los siguientes supuestos:
\begin{lnum}
    \item El esfuerzo de uno de los factores de producción es endógeno y constituye un mal para éste. Como postula la ley psicológica de YERKES-DODSON, los individuos que están expuestos a baja presión no pondrán mucho esfuerzo en sus acciones. Según crece la presión, la situación cambia hasta alcanzar el punto en que demasiada presión resulta ineficiente;
    \item Contratos incompletos (i.e. no pueden prever todas las contingencias). Los contratos de empleo definen la compensación económica de los trabajadores pero no otras preocupaciones como el esfuerzo, la carga de trabajo o tareas específicas que no están definidas, lo que crea un vacío en los contratos;
    \item Inercia. Las fuerzas que son externas a la compañía pueden ejercer presión para hacer que la empresa sea más competitiva. Solo con estas fuerzas el esfuerzo aumentará; y,
    \item Prudencia (información asimétrica). Se asume que los trabajadores no mostrarán su pleno potencial y que al mismo tiempo, los empleadores no pagarán el máximo salario que pueden ofrecer.
\end{lnum}

Por todo lo anterior, los trabajadores (que son los que deciden el nivel de esfuerzo a llevar a cabo), escogerán por lo general un grado de esfuerzo inferior al que maximiza beneficios. Así, el comportamiento optimizador (i.e. producción que minimiza el coste medio) es una excepción, produciéndose generalmente una ineficiencia X.

Es muy interesante como esta ineficiencia de X planteada por LEIBENSTEIN se desarrollo más en profundidad en la Economía de la Información [\authorfont{Ver Tema 3.A.13}] Además, el grado de competencia juega un papel fundamental en la teoría de LEIBENSTEIN:
\begin{lnum}
    \item Con competencia alta, los trabajadores, en aras de asegurar la supervivencia de la empresa (y conservar, por tanto, su puesto de trabajo), se esforzarán y mejorarán su productividad marginal;
    \item Con competencia baja, los trabajadores tendrán incentivos a relajares. Así, un aumento de los precios del producto no tiene por qué traducirse en un mayor beneficio, ya que si los trabajadores perciben que ese aumento del precio de venta se debe a un mark-up posibilitado por una menor competencia, entonces el beneficio podría incluso disminuir al aumentar los costes por la ineficiencia X; y,
    \item Sin embargo, pruebas empíricas y otros modelos sobre la relevancia de las ineficiencias X tienen resultados mixtos por lo que no se puede llegar a ninguna conclusión. 
\end{lnum}



\clearpage
\section{Aplicaciones empíricas: Métodos de estimación}

\subsection{Análisis envolvente de datos: DEA}
Desarrollado por CHARNES, COOPER y RHODES en 1978, el análisis envolvente de datos surgió como respuesta a las limitaciones de los métodos tradicionales de medición de la eficiencia. Se trata de un método no-paramétrico utilizado en investigación operativa para medir o evaluar la eficiencia de las unidades de toma de decisiones (DMU’s), como empresas, servicios públicos o cualquier entidad que convierta inputs en outputs. El DEA es particularmente útil en situaciones en las que intervienen múltiples inputs y/o outputs, lo que permite una evaluación integral de la eficiencia relativa. Al emplear métodos de programación lineal, el DEA construye una frontera de rendimiento óptimo, lo que permite la comparación de entidades similares.

Aunque aparentemente conduce a la creación de una Frontera de Posibilidades de Producción, en puridad, al ser un proceso de benchmarking la evaluación a través del DEA conduce a la creación de una “Frontera de Mejores Prácticas” sobre la que se evalúa el rendimiento de los DMU’s. 

\subsubsection{Modelo básico}
Se busca asignar un valor de eficiencia a cada DMU en base a los inputs y outputs observados en éste, sujeto a la restricción impuesta por el resto de DMU’s y por otros de consistencia. En concreto, se maximiza el siguiente DMU:

\eqblock{
\begin{aligned}
    &\boxed{\begin{aligned}
        &\max\qquad \theta_j=\frac{\sum_{m=1}^My_m^ju_m^j}{\sum_{n=1}^Nx_n^jv_n^j}\\[1em]
        &\text{s.a.}\qquad 
        \begin{cases}
            \frac{\sum_{m=1}^My_m^ju_m^j}{\sum_{n=1}^Nx_n^jv_n^j}\leq1\qquad \forall k=1,\dots,K\\[1em]
            \text{\scriptsize
            $\begin{aligned}
                &0\ll(y_1^j,y_2^j,\dots, y_M^j)\\
                &0\ll(u_1^j,u_2^j\dots, u_M^j)\\
                &0\ll(x_1^j.x_2^j,\dots,x_N^j)\\
                &0\ll(v_1^j,v_2^j,\dots,v_N^j)
            \end{aligned}$}
        \end{cases}
    \end{aligned}}\\[2em]
    &\left\{\begin{aligned}
        &(y_1^j,y_2^j,\dots, y_M^j)\equiv &&\text{Outputs del proceso productivo}\\
        &(u_1^j,u_2^j\dots, u_M^j)\equiv &&\text{Pesos relativos de cada output producido}\\
        &(x_1^j.x_2^j,\dots,x_N^j)\equiv &&\text{Inputs del proceso productivo}\\
        &(v_1^j,v_2^j,\dots,v_N^j)\equiv &&\text{Pesos relativos de cada input necesario}
    \end{aligned}\right.
\end{aligned}
}
 
Sujeto a la restricción de no-negatividad de inputs, outputs y pesos relativos inferiores o iguales a uno (y mayores o iguales a cero), y una adicional (primera): que la eficiencia de cada sistema tiene que ser menor o igual a uno, en todo caso. Al tratarse de un problema de programación lineal no podremos utilizar cocientes, entonces deberemos reformular las ecuaciones tal que así:

$$\begin{aligned}
    &\max\qquad \theta_j\sum_{m=1}^My_m^ju_m^j\\[0.5em]
    \text{s.a.}\qquad
    &\left\{\begin{aligned}
        &\begin{matrix}
            \sum_{m=1}^My_m^1u_m^1-\sum_{n=1}^Nx_n^1v_n^1\leq0\\[0.5em]
            \vdots\;(\forall j=1,\dots,J)\\[0.5em]
            \sum_{m=1}^My_m^Ju_m^J-\sum_{n=1}^Nx_n^Jv_n^J\leq0\\[0.5em]
        \end{matrix}\\[0.5em]
        &\sum_{m=1}^Mx_m^jv_m^j=1\\[0.5em]
        &0\leq v_n,u_n
    \end{aligned}\right.
\end{aligned}$$

Con estas restricciones obtendremos un ránking de eficiencia (comparativa) y una relación de inputs que defina la frontera eficiente. De forma que:

\imagenfit{Frontera_DEA.png}{Frontera eficiente (\textit{benchmarking}): Análisis Envolvente de Datos (DEA)}{Elaboración propia}
 
Para medir la Ratio de ineficiencia deberíamos hacer ($1 – \frac{\over{(0,A')}}{\over{(0,A)}}$). Además de este resultado, podemos obtener también lo siguiente:
\begin{lnum}
    \item Ponderaciones, o precios de los inputs y outputs (soluciones v y u): indica el peso asignado a cada variable output e input en la formación del ratio de eficiencia;
    \item Reducción proporcional de los inputs, para el caso de los modelos orientados a los inputs. Esto equivaldría a la minimización de la distancia radial desde la unidad ineficiente a la frontera. En estos se busca la eficiencia técnica de la producción. En el modelo orientado a los inputs nos estaríamos moviendo hacía dentro de la nube de datos (hacia el origen); y,
    \item Aumento proporcional de los outputs. El incremento proporcional que debe experimentar la unidad ineficiente en su eficiencia para llegar a la frontera. En el modelo orientado a los outputs nos estaríamos moviendo hacía afuera de la nube de datos (expansión del output).\footnote{%
        La elección del problema de minimización o maximización depende del criterio del analista y de su conocimiento del sector objeto del análisis. En determinados problemas, las empresas no tendrán libertad a la hora de fijar la cantidad de output por lo que será mejor un enfoque de minimización de inputs. Por el contrario, en otros contextos la cantidad de inputs puede venir establecida por una regulación y el operador se concentrará entonces en maximizar su output.
    }
\end{lnum}
	 
\subsubsection{Extensiones}
Sin entrar en materia, debemos destacar que existen también diferentes extensiones o variaciones del modelo básico de DEA planteado. En particular, respecto a:
\begin{lnum}
    \item Rendimientos a escala. El analista podrá optar por considerar si nos encontramos en el largo plazo, supuesto altamente restrictivo pero que nos permite calcular la eficiencia técnica y asignativa en presencia de rendimientsos variables a escala;
    \item Análisis estático o dinámico. En caso de contar con un panel quilibrado de datos podremos, además de analizar el DEA para un determinado periodo, incorporar una opción dinámica que aproveche esta dimensión temporal, permitiéndonos calcular un índice de productividad total de los factores (Malmquist Index, comparación bilateral) para cada periodo y empresa.\footnote{%
        Ver: \href{https://en.wikipedia.org/wiki/Malmquist_index}{Malmquist index (Wikipedia.org)}
    }
\end{lnum}

\subsubsection{Limitaciones}
A pesar de sus ventajas, el análisis envolvente de datos tiene sus propias limitaciones. Una desventaja importante es su sensibilidad a outliers, lo que puede distorsionar los puntajes de eficiencia y llevar a conclusiones erróneas. Además, DEA no tiene en cuenta la naturaleza estocástica (no-paramétrico) de los datos, lo que significa que supone que todas las desviaciones de la frontera se deben a la ineficiencia y no a la variación aleatoria. Esta limitación puede afectar la confiabilidad de los resultados, en particular en conjuntos de datos con alta variabilidad. Algunas alternativas propuestas para mejorar alguno de sus inconvenientes o potenciar sus ventajas son: el DEA Estocástico  o el Stochastic Frontier Analysis.\footnote{%
    Ver: \href{https://www.sciencedirect.com/science/article/abs/pii/S0377221715007134}{\textit{Stochastic Data Envelopment Analysis—A review} (OLESSEN, 2016)}
}

\subsection{Stochastic Frontier Analysis}
Simultaneamente introducido por AIGNER, LOVELL y SCHMIDT (1977) and MEEUSEN and VAN DEN BROECK (1977), el análisis fontera estocástica es uno de los métodos más utilizados para analizar y medir la eficiencia. Se trata de una sofisticación de los modelos más tradicionales de frontera de producción. 

\subsubsection{Modelo}
Podemos definir el modelo frontera de producción estocástica como:

\eqblock{
\begin{aligned}
    \boxed{y_i=f(\vec{x_i},\vec\beta)\cdot e^{v_i}\cdot E_i}\\[2em]
    \begin{aligned}
        &y,x\equiv \text{Inputs y outputs del productor $i$}\\[0.5em]
        &\underset{\text{\scriptsize Debe venir explícitamente definida, por lo que es una limitación importante}}{f(\vec{x_i})\equiv\text{Función de Producción del productor $i$}}\\[0.5em]
        &\underset{\text{\scriptsize Se trata de los parámetros a estimar (parámetros del modelo)}}{\vec{\beta}\equiv\text{Coeficientes de producción}}\\[0.5em]
        &v_i\equiv\text{Componente aleatorio para productor $i$}\\[0.5em]
        &E_i\equiv \text{Término de eficiencia ténica para productor $i$}
    \end{aligned}
\end{aligned}
}{Análisis por frontera estocástica (verosimilitud): Formulación del modelo general. Aplicaciones empíricas: Métodos de estimación}

El componente aleatorio puede eliminarse del modelo, igual que el componente de eficiencia técnica puede expresarse en forma exponencial (negativa, ya que estará acotado entre 0 y 1). De hecho, si suponemos que la forma funcional es del tipo COBB-DOUGLAS, formularemos la expresión en su forma logarítmica tal que:

\eqblock{
\begin{aligned}
    &\ln{y_i}=\overbrace{\beta_0+\sum_{n=1}^N\beta_nx_{i,n}+v_i}^{\text{Frontera estocástica}}-u_i\qquad \text{y,}\qquad \varepsilon_i=v_i-u_i\\[0.5em]
    &\Longrightarrow\qquad \boxed{\ln{y_i}=\beta_0+\sum_{n=1}^N\beta_nx_{i,n}+\varepsilon_i}
\end{aligned}
}{Análisis por frontera estocástica (verosimilitud): Formulación para forma funcional COBB-DOUGLAs. Aplicaciones empíricas: Métodos de estimación}

Ahora bien, esto es solo el inicio del problema, es decir, no pretendemos estimar los parámetros mediante MCO’s sino que nuestra intención será encontrar la distribución del error compuesto, ajustando los parámetros, de tal forma que se ajuste lo máximo posible a los datos obtenidos, siempre asumiendo o teniendo:
\begin{lnum}
    \item $v$ es una variable de ruido aleatorio (captura errores de medición, clima, suerte, etc.) sobre la que se asume (generalmente) una distribución $v\sim N(0,\sigma_v)$; y,
    \item $u$ corresponde a la ineficiencia técnica (captura cuánto se desvía la empresa de la frontera estocástica), para esta se asumirá, por norma general, una distribución $u\sim N(0,\sigma_u)\;(+)$ (solo la parte positiva, pues la ineficiencia no puede ser negativa).
\end{lnum}
	
Para conseguirlo, el método más utilizado es el de Máxima Verosimiliud  que, usualmente, es el que proporciona resultados más eficientes si la suposición de independencia tanto de los efectos como de los regresores es plausible y tiene sentido económico. 

\begin{specialblock}{Método de Máxima Verosimilitud}
    La función de verosimilitud (o, simplemente, verosimilitud) es una función de los parámetros de un modelo estadístico (por MCO’s por ejemplo) que permite realizar inferencias acerca de su valor a partir de un conjunto de observaciones. En cierto sentido, la verosimilitud es una versión inversa de la probabilidad condicional.
    
    Vamos a seguir el método de Máxima Verosimilitud por lo siguiente: en el modelo planteado lo que nos interesa obtener es el término de ineficiencia, que es lo verdaderamente interesante. Ahora bien, la introducción de ese término altera la distribución “generalmente aceptada” del término de error (la Normal) puesto que, al ser ahora un error compuesto, el error ya no es simétrico por ser la diferencia entre un ruido normal y un término positivo (la ineficiencia): su distribución será ahora asimétrica a la izquierda, entonces, el MCO’s ya no servirá, necesitaremos un método que ajuste esa distribución asimétrica $\to$ La Máxima Verosimilitud.
\end{specialblock}

El Método de Máxima Verosimilitud no estima directamente una recta por MCO’s, sino que busca los parámetros (parámetros de la función y varianzas tanto del error ruido y del error ineficiencia, cada uno por separado) que maximizan la probabilidad de observar los datos de nuestra muestra:

$$\begin{aligned}
    L(\beta,\sigma_v,\sigma_u)=\Pi_{i=1}^nf_\varepsilon(\varepsilon_i;\;\sigma_v,\sigma_u)
\end{aligned}$$
 
Donde $f_\varepsilon$ es la función de densidad del error compuesto, derivada de las distribuciones asumidas de los dos componentes. Mediante integración, obtendremos una función de densidad no simétrica que será la función de verosimilitud individual para cada observación. 

Como las observaciones son independientes, la función de verosimilitud total se obtendrá de:

$$\begin{aligned}
    L(\beta,\sigma_v,\sigma_u)=\Pi_{i=1}^nf(\varepsilon_i)
\end{aligned}$$
 
Y, si aplicamos logarítmos, trabajaremos con el sumatorio de funciones individuales. A partir de aquí, el método de Máxima Verosimilitud ajusta los parámetros mediante ciclos iterativos hasta que la log-verosimilitud alcance su máximo (lo que implicaría que la probabilidad de obtener los resultados de nuestra muestra dados esos parámetros es la más alta, sería como una reversión del proceso de MCO’s). El modelo final nos entregará:
\begin{lnum}
    \item Las betas, elasticidades de los inputs;
    \item Las varianzas de ineficiencia y ruida; y,
    \item Y las medidas individuales de eficiencia, esto es:
    $$\mathbb E[u_i\;|\;\varepsilon_i]$$
    Esto es lo que realmente nos interesa, puesto que, para calcular la eficiencia técnica de cada empresa solo tendremos que hacer laexponencial a la medida anterior. Y, además, podremos calcular el de cualquier otra que queramos ya que esto consiste en: la ineficiencia esperada de la empresa dado el ruido observado en la regresión. 
\end{lnum}
	
\begin{specialblock}{\textbf{Ejemplos.-} Estimaciones en sectores económicos}
    \href{https://www.sciencedirect.com/science/article/abs/pii/S0140988310000496}{\textit{Measurement of productive efficiency with frontier methods: A case study for wind farms}, IGLESIAS, et. al., 2010}

    \href{}{\textit{The impact of domestic and foreign R\&D on agricultural productivity in sub-Saharan Africa} ADETUTU, et. al. (2019).}

    \href{https://doi.org/10.1080/15140326.2023.2178798}{\textit{Technical efficiency of U.S. Western Great Plains wheat farms using stochastic frontier analysis} JI, VITALE and ADAM, 2023}

    \href{https://www.cambridge.org/core/journals/agricultural-and-resource-economics-review/article/efficiency-analysis-of-developing-country-agriculture-a-review-of-the-frontier-function-literature/9F5FB5B35C8DE12797588E3B5B41CF44}{\textit{Efficiency Analysis of Developing Country Agriculture: A review of the Frontier Function Literature} BRAVO-URETA and PINHEIRO, 2016}

    \href{https://so01.tci-thaijo.org/index.php/AEJ/article/view/242756/166696}{\textit{A comparison of DEA and SFA approaches: Application to the US Non-Life Insurance Market} FERNANDEZ, QUIROGA-GARCIA and MANZANO-PEREZ, 2020}
\end{specialblock}















\clearpage
\section{Conclusión} 

\subsection{Extensiones y relación con otras partes del Temario}


\subsection{Opinión}



\subsubsection{Idea final (Salida o cierra)}


\clearpage
\pagenumbering{gobble}
\MyBibliography
\MyBibItem{Autor}{Título}{Año}{DOI -- LINK}




%ANEXOS
\clearpage
\anexo{\textbf{Preguntas Test}}
%\input{\TemaCode/questions.tex} 




\end{document}